% Options for packages loaded elsewhere
\PassOptionsToPackage{unicode}{hyperref}
\PassOptionsToPackage{hyphens}{url}
\documentclass[
]{article}
\usepackage{xcolor}
\usepackage[margin=1in]{geometry}
\usepackage{amsmath,amssymb}
\setcounter{secnumdepth}{-\maxdimen} % remove section numbering
\usepackage{iftex}
\ifPDFTeX
  \usepackage[T1]{fontenc}
  \usepackage[utf8]{inputenc}
  \usepackage{textcomp} % provide euro and other symbols
\else % if luatex or xetex
  \usepackage{unicode-math} % this also loads fontspec
  \defaultfontfeatures{Scale=MatchLowercase}
  \defaultfontfeatures[\rmfamily]{Ligatures=TeX,Scale=1}
\fi
\usepackage{lmodern}
\ifPDFTeX\else
  % xetex/luatex font selection
\fi
% Use upquote if available, for straight quotes in verbatim environments
\IfFileExists{upquote.sty}{\usepackage{upquote}}{}
\IfFileExists{microtype.sty}{% use microtype if available
  \usepackage[]{microtype}
  \UseMicrotypeSet[protrusion]{basicmath} % disable protrusion for tt fonts
}{}
\makeatletter
\@ifundefined{KOMAClassName}{% if non-KOMA class
  \IfFileExists{parskip.sty}{%
    \usepackage{parskip}
  }{% else
    \setlength{\parindent}{0pt}
    \setlength{\parskip}{6pt plus 2pt minus 1pt}}
}{% if KOMA class
  \KOMAoptions{parskip=half}}
\makeatother
\usepackage{color}
\usepackage{fancyvrb}
\newcommand{\VerbBar}{|}
\newcommand{\VERB}{\Verb[commandchars=\\\{\}]}
\DefineVerbatimEnvironment{Highlighting}{Verbatim}{commandchars=\\\{\}}
% Add ',fontsize=\small' for more characters per line
\usepackage{framed}
\definecolor{shadecolor}{RGB}{248,248,248}
\newenvironment{Shaded}{\begin{snugshade}}{\end{snugshade}}
\newcommand{\AlertTok}[1]{\textcolor[rgb]{0.94,0.16,0.16}{#1}}
\newcommand{\AnnotationTok}[1]{\textcolor[rgb]{0.56,0.35,0.01}{\textbf{\textit{#1}}}}
\newcommand{\AttributeTok}[1]{\textcolor[rgb]{0.13,0.29,0.53}{#1}}
\newcommand{\BaseNTok}[1]{\textcolor[rgb]{0.00,0.00,0.81}{#1}}
\newcommand{\BuiltInTok}[1]{#1}
\newcommand{\CharTok}[1]{\textcolor[rgb]{0.31,0.60,0.02}{#1}}
\newcommand{\CommentTok}[1]{\textcolor[rgb]{0.56,0.35,0.01}{\textit{#1}}}
\newcommand{\CommentVarTok}[1]{\textcolor[rgb]{0.56,0.35,0.01}{\textbf{\textit{#1}}}}
\newcommand{\ConstantTok}[1]{\textcolor[rgb]{0.56,0.35,0.01}{#1}}
\newcommand{\ControlFlowTok}[1]{\textcolor[rgb]{0.13,0.29,0.53}{\textbf{#1}}}
\newcommand{\DataTypeTok}[1]{\textcolor[rgb]{0.13,0.29,0.53}{#1}}
\newcommand{\DecValTok}[1]{\textcolor[rgb]{0.00,0.00,0.81}{#1}}
\newcommand{\DocumentationTok}[1]{\textcolor[rgb]{0.56,0.35,0.01}{\textbf{\textit{#1}}}}
\newcommand{\ErrorTok}[1]{\textcolor[rgb]{0.64,0.00,0.00}{\textbf{#1}}}
\newcommand{\ExtensionTok}[1]{#1}
\newcommand{\FloatTok}[1]{\textcolor[rgb]{0.00,0.00,0.81}{#1}}
\newcommand{\FunctionTok}[1]{\textcolor[rgb]{0.13,0.29,0.53}{\textbf{#1}}}
\newcommand{\ImportTok}[1]{#1}
\newcommand{\InformationTok}[1]{\textcolor[rgb]{0.56,0.35,0.01}{\textbf{\textit{#1}}}}
\newcommand{\KeywordTok}[1]{\textcolor[rgb]{0.13,0.29,0.53}{\textbf{#1}}}
\newcommand{\NormalTok}[1]{#1}
\newcommand{\OperatorTok}[1]{\textcolor[rgb]{0.81,0.36,0.00}{\textbf{#1}}}
\newcommand{\OtherTok}[1]{\textcolor[rgb]{0.56,0.35,0.01}{#1}}
\newcommand{\PreprocessorTok}[1]{\textcolor[rgb]{0.56,0.35,0.01}{\textit{#1}}}
\newcommand{\RegionMarkerTok}[1]{#1}
\newcommand{\SpecialCharTok}[1]{\textcolor[rgb]{0.81,0.36,0.00}{\textbf{#1}}}
\newcommand{\SpecialStringTok}[1]{\textcolor[rgb]{0.31,0.60,0.02}{#1}}
\newcommand{\StringTok}[1]{\textcolor[rgb]{0.31,0.60,0.02}{#1}}
\newcommand{\VariableTok}[1]{\textcolor[rgb]{0.00,0.00,0.00}{#1}}
\newcommand{\VerbatimStringTok}[1]{\textcolor[rgb]{0.31,0.60,0.02}{#1}}
\newcommand{\WarningTok}[1]{\textcolor[rgb]{0.56,0.35,0.01}{\textbf{\textit{#1}}}}
\usepackage{longtable,booktabs,array}
\newcounter{none} % for unnumbered tables
\usepackage{calc} % for calculating minipage widths
% Correct order of tables after \paragraph or \subparagraph
\usepackage{etoolbox}
\makeatletter
\patchcmd\longtable{\par}{\if@noskipsec\mbox{}\fi\par}{}{}
\makeatother
% Allow footnotes in longtable head/foot
\IfFileExists{footnotehyper.sty}{\usepackage{footnotehyper}}{\usepackage{footnote}}
\makesavenoteenv{longtable}
\usepackage{graphicx}
\makeatletter
\newsavebox\pandoc@box
\newcommand*\pandocbounded[1]{% scales image to fit in text height/width
  \sbox\pandoc@box{#1}%
  \Gscale@div\@tempa{\textheight}{\dimexpr\ht\pandoc@box+\dp\pandoc@box\relax}%
  \Gscale@div\@tempb{\linewidth}{\wd\pandoc@box}%
  \ifdim\@tempb\p@<\@tempa\p@\let\@tempa\@tempb\fi% select the smaller of both
  \ifdim\@tempa\p@<\p@\scalebox{\@tempa}{\usebox\pandoc@box}%
  \else\usebox{\pandoc@box}%
  \fi%
}
% Set default figure placement to htbp
\def\fps@figure{htbp}
\makeatother
\setlength{\emergencystretch}{3em} % prevent overfull lines
\providecommand{\tightlist}{%
  \setlength{\itemsep}{0pt}\setlength{\parskip}{0pt}}
\usepackage{bookmark}
\IfFileExists{xurl.sty}{\usepackage{xurl}}{} % add URL line breaks if available
\urlstyle{same}
\hypersetup{
  pdftitle={Exploración de datos · SDA Lab},
  hidelinks,
  pdfcreator={LaTeX via pandoc}}

\title{Exploración de datos · SDA Lab}
\usepackage{etoolbox}
\makeatletter
\providecommand{\subtitle}[1]{% add subtitle to \maketitle
  \apptocmd{\@title}{\par {\large #1 \par}}{}{}
}
\makeatother
\subtitle{ori}
\author{}
\date{\vspace{-2.5em}2026-09-12 07:50}

\begin{document}
\maketitle

Cuaderno generado desde SDA Lab. Cada sección corresponde a un panel
marcado con la casilla \textbf{Añadir}: trae su texto explicativo, los
parámetros con que se produjo y el código R que lo redibuja sobre
\texttt{datos}.

\section{Los datos}\label{los-datos}

{\def\LTcaptype{none} % do not increment counter
\begin{longtable}[]{@{}ll@{}}
\toprule\noalign{}
Campo & Valor \\
\midrule\noalign{}
\endhead
\bottomrule\noalign{}
\endlastfoot
Nombre & ori \\
Fuente & ori \\
Filas al cargar & 4543 \\
Columnas al cargar & 18 \\
Filas ahora & 118 \\
Columnas ahora & 18 \\
\end{longtable}
}

\begin{Shaded}
\begin{Highlighting}[]
\CommentTok{\# ORI.csv es la base del Taller 01 (IDEAM). Dejalo al lado del cuaderno o}
\CommentTok{\# ajustá \textasciigrave{}ruta\textasciigrave{}: se prueban los lugares habituales, en orden.}
\NormalTok{ruta }\OtherTok{\textless{}{-}} \FunctionTok{Filter}\NormalTok{(file.exists, }\FunctionTok{c}\NormalTok{(}\StringTok{"ORI.csv"}\NormalTok{, }\StringTok{"data/ORI.csv"}\NormalTok{,}
                              \StringTok{"../../../data/ORI.csv"}\NormalTok{))[}\DecValTok{1}\NormalTok{]}
\CommentTok{\# La copia que reparte el curso viene en ISO{-}8859{-}1 y la de data/ del repo}
\CommentTok{\# en UTF{-}8. Se detecta en vez de suponer: suponer es lo que rompe las}
\CommentTok{\# tildes de los municipios (ACACÍAS {-}\textgreater{} ACACÃAS).}
\NormalTok{codificacion }\OtherTok{\textless{}{-}} \ControlFlowTok{if}\NormalTok{ (}\FunctionTok{all}\NormalTok{(}\FunctionTok{validUTF8}\NormalTok{(}\FunctionTok{readLines}\NormalTok{(ruta, }\AttributeTok{warn =} \ConstantTok{FALSE}\NormalTok{,}
                                            \AttributeTok{encoding =} \StringTok{"bytes"}\NormalTok{))))}
  \StringTok{"UTF{-}8"} \ControlFlowTok{else} \StringTok{"latin1"}
\NormalTok{datos }\OtherTok{\textless{}{-}} \FunctionTok{read.csv}\NormalTok{(ruta, }\AttributeTok{sep =} \StringTok{";"}\NormalTok{, }\AttributeTok{dec =} \StringTok{"."}\NormalTok{, }\AttributeTok{fileEncoding =}\NormalTok{ codificacion,}
                  \AttributeTok{stringsAsFactors =} \ConstantTok{FALSE}\NormalTok{)}
\FunctionTok{names}\NormalTok{(datos) }\OtherTok{\textless{}{-}} \FunctionTok{trimws}\NormalTok{(}\FunctionTok{names}\NormalTok{(datos))   }\CommentTok{\# el encabezado trae "Pronostico "}
\FunctionTok{cat}\NormalTok{(}\StringTok{"Al cargar:"}\NormalTok{, }\FunctionTok{nrow}\NormalTok{(datos), }\StringTok{"filas x"}\NormalTok{, }\FunctionTok{ncol}\NormalTok{(datos), }\StringTok{"columnas}\SpecialCharTok{\textbackslash{}n}\StringTok{"}\NormalTok{)}
\end{Highlighting}
\end{Shaded}

\begin{verbatim}
## Al cargar: 4543 filas x 18 columnas
\end{verbatim}

\section{Preparación}\label{preparaciuxf3n}

\begin{Shaded}
\begin{Highlighting}[]
\NormalTok{datos }\OtherTok{\textless{}{-}}\NormalTok{ datos[datos}\SpecialCharTok{$}\NormalTok{Hora }\SpecialCharTok{\%in\%} \FunctionTok{c}\NormalTok{(}\StringTok{"12:00"}\NormalTok{), ]}
\FunctionTok{cat}\NormalTok{(}\StringTok{"Después de la preparación:"}\NormalTok{, }\FunctionTok{nrow}\NormalTok{(datos), }\StringTok{"filas}\SpecialCharTok{\textbackslash{}n}\StringTok{"}\NormalTok{)}
\end{Highlighting}
\end{Shaded}

\begin{verbatim}
## Después de la preparación: 118 filas
\end{verbatim}

\section{Análisis}\label{anuxe1lisis}

\subsection{Vista previa del dataset}\label{vista-previa-del-dataset}

Panel \texttt{f1.fuente.vista\_previa} · marcado a las 07:25:51

\subsubsection{Para qué sirve}\label{para-quuxe9-sirve}

Atrapar el desastre antes de gastar una hora en él. Es la única card que
se mira sin haber decidido nada todavía, y la que dice si el archivo se
cargó como creías o hay que volver a leerlo con otras opciones.

\subsubsection{Qué muestra}\label{quuxe9-muestra}

Las primeras filas de lo que se acaba de cargar, con el pie que dice
cuántas se muestran y cuántas hay. Es la comprobación más barata que
existe y la que más errores atrapa: si la tabla se ve mal acá, todo lo
que venga después está mal.

\subsubsection{Qué buscar}\label{quuxe9-buscar}

\begin{itemize}
\tightlist
\item
  \textbf{Columnas de más o de menos}: casi siempre es el separador
  equivocado en un CSV. Si todo el archivo entró en una sola columna,
  probá \texttt{;} o tabulador.
\item
  \textbf{Números leídos como texto}: el separador decimal. Con coma
  decimal y punto como separador de miles, \texttt{1.234,5} se convierte
  en texto y ninguna operación numérica va a funcionar.
\item
  \textbf{Tildes rotas}: es la codificación. UTF-8 y Latin-1 se
  confunden fácil.
\item
  \textbf{Faltantes disfrazados}: \texttt{.}, \texttt{-99},
  \texttt{N/A}, celdas vacías. Si no se declaran como faltantes, entran
  al análisis como valores reales.
\item
  \textbf{La primera fila}: ¿es un encabezado o ya es un dato?
\end{itemize}

\subsubsection{Cuándo engaña}\label{cuuxe1ndo-engauxf1a}

\textbf{Ver diez filas correctas no dice nada de las 35.000 restantes.}
El problema típico ---una columna que cambia de formato a mitad del
archivo--- no aparece arriba. La subsección de Calidad es la que mira el
archivo entero.

\textbf{Un dataset del curso ya cargado también puede estar
transformado.} Si en el pivot país×año ves ceros, muchos no son ceros
observados: son pares sin dato rellenados en la agregación. El panel lo
avisa al cargar y conviene creerle.

\textbf{El orden de las filas puede no ser inocente.} Si el archivo
viene ordenado por alguna variable, cualquier partición sin barajar
hereda ese orden. Por eso la subsección de Partición baraja siempre y
con semilla.

Estado al marcarlo:

{\def\LTcaptype{none} % do not increment counter
\begin{longtable}[]{@{}
  >{\raggedright\arraybackslash}p{(\linewidth - 34\tabcolsep) * \real{0.0556}}
  >{\raggedright\arraybackslash}p{(\linewidth - 34\tabcolsep) * \real{0.0556}}
  >{\raggedright\arraybackslash}p{(\linewidth - 34\tabcolsep) * \real{0.0556}}
  >{\raggedright\arraybackslash}p{(\linewidth - 34\tabcolsep) * \real{0.0556}}
  >{\raggedright\arraybackslash}p{(\linewidth - 34\tabcolsep) * \real{0.0556}}
  >{\raggedright\arraybackslash}p{(\linewidth - 34\tabcolsep) * \real{0.0556}}
  >{\raggedright\arraybackslash}p{(\linewidth - 34\tabcolsep) * \real{0.0556}}
  >{\raggedright\arraybackslash}p{(\linewidth - 34\tabcolsep) * \real{0.0556}}
  >{\raggedright\arraybackslash}p{(\linewidth - 34\tabcolsep) * \real{0.0556}}
  >{\raggedright\arraybackslash}p{(\linewidth - 34\tabcolsep) * \real{0.0556}}
  >{\raggedright\arraybackslash}p{(\linewidth - 34\tabcolsep) * \real{0.0556}}
  >{\raggedright\arraybackslash}p{(\linewidth - 34\tabcolsep) * \real{0.0556}}
  >{\raggedright\arraybackslash}p{(\linewidth - 34\tabcolsep) * \real{0.0556}}
  >{\raggedright\arraybackslash}p{(\linewidth - 34\tabcolsep) * \real{0.0556}}
  >{\raggedright\arraybackslash}p{(\linewidth - 34\tabcolsep) * \real{0.0556}}
  >{\raggedright\arraybackslash}p{(\linewidth - 34\tabcolsep) * \real{0.0556}}
  >{\raggedright\arraybackslash}p{(\linewidth - 34\tabcolsep) * \real{0.0556}}
  >{\raggedright\arraybackslash}p{(\linewidth - 34\tabcolsep) * \real{0.0556}}@{}}
\toprule\noalign{}
\begin{minipage}[b]{\linewidth}\raggedright
Cod\_Div
\end{minipage} & \begin{minipage}[b]{\linewidth}\raggedright
Latitud
\end{minipage} & \begin{minipage}[b]{\linewidth}\raggedright
Longitud
\end{minipage} & \begin{minipage}[b]{\linewidth}\raggedright
Region
\end{minipage} & \begin{minipage}[b]{\linewidth}\raggedright
Departamento
\end{minipage} & \begin{minipage}[b]{\linewidth}\raggedright
Municipio
\end{minipage} & \begin{minipage}[b]{\linewidth}\raggedright
Fecha
\end{minipage} & \begin{minipage}[b]{\linewidth}\raggedright
Hora
\end{minipage} & \begin{minipage}[b]{\linewidth}\raggedright
Temperatura
\end{minipage} & \begin{minipage}[b]{\linewidth}\raggedright
Velocidad\_del\_Viento
\end{minipage} & \begin{minipage}[b]{\linewidth}\raggedright
Direccion\_del\_Viento
\end{minipage} & \begin{minipage}[b]{\linewidth}\raggedright
Presion
\end{minipage} & \begin{minipage}[b]{\linewidth}\raggedright
Punto\_de\_Rocio
\end{minipage} & \begin{minipage}[b]{\linewidth}\raggedright
Cobertura\_total\_nubosa
\end{minipage} & \begin{minipage}[b]{\linewidth}\raggedright
Precipitacion\_mm\_h
\end{minipage} & \begin{minipage}[b]{\linewidth}\raggedright
Probabilidad\_de\_Tormenta
\end{minipage} & \begin{minipage}[b]{\linewidth}\raggedright
Humedad
\end{minipage} & \begin{minipage}[b]{\linewidth}\raggedright
Pronostico
\end{minipage} \\
\midrule\noalign{}
\endhead
\bottomrule\noalign{}
\endlastfoot
50006000 & 3.99012 & -73.76594 & ORI & META & ACACÍAS & 26/08/2022 &
12:00 & 26.5 & 1.2 & 179.8 & 1014.2 & 18.3 & 70 & 0 & 0 & 60.7 &
Parcialmente Nublado \\
50006000 & 3.99012 & -73.76594 & ORI & META & ACACÍAS & 27/08/2022 &
12:00 & 27.3 & 1.7 & 113.7 & 1013.9 & 21 & 80 & 0 & 0 & 68.7 &
Nublado \\
85010000 & 5.17269 & -72.5463 & ORI & CASANARE & AGUAZUL & 26/08/2022 &
12:00 & 30.2 & 2.1 & 176.7 & 1013 & 21.8 & 49.8 & 0 & 0 & 60.9 &
Parcialmente Nublado \\
85010000 & 5.17269 & -72.5463 & ORI & CASANARE & AGUAZUL & 27/08/2022 &
12:00 & 30.4 & 1.7 & 200 & 1013.2 & 24.4 & 63.2 & 0 & 0 & 70.4 &
Parcialmente Nublado \\
81001000 & 7.073902 & -70.74801 & ORI & ARAUCA & ARAUCA & 26/08/2022 &
12:00 & 30.4 & 2 & 150 & 1012 & 22.5 & 40 & 0 & 0 & 62.7 & Parcialmente
Nublado \\
81001000 & 7.073902 & -70.74801 & ORI & ARAUCA & ARAUCA & 27/08/2022 &
12:00 & 29 & 0.7 & 187.1 & 1011.6 & 22.3 & 80 & 0 & 0 & 67.5 &
Nublado \\
81065000 & 7.026981 & -71.42675 & ORI & ARAUCA & ARAUQUITA & 26/08/2022
& 12:00 & 29.7 & 2.1 & 200.7 & 1011.6 & 18.3 & 50 & 0 & 0 & 50.3 &
Parcialmente Nublado \\
81065000 & 7.026981 & -71.42675 & ORI & ARAUCA & ARAUQUITA & 27/08/2022
& 12:00 & 28.8 & 0.5 & 356.5 & 1011.9 & 24.5 & 80 & 0.8 & 0 & 77.6 &
Nublado - Lluvia \\
50110000 & 4.567202 & -72.96328 & ORI & META & BARRANCA DE UPÍA &
26/08/2022 & 12:00 & 31 & 2 & 200 & 1012.2 & 21.4 & 50 & 0 & 0 & 56.7 &
Parcialmente Nublado \\
50110000 & 4.567202 & -72.96328 & ORI & META & BARRANCA DE UPÍA &
27/08/2022 & 12:00 & 31.5 & 0.7 & 210.5 & 1012.8 & 28.7 & 80 & 0 & 0 &
85 & Nublado \\
\end{longtable}
}

\begin{Shaded}
\begin{Highlighting}[]
\CommentTok{\# f1.fuente.vista\_previa}
\FunctionTok{cat}\NormalTok{(}\StringTok{"Unidades estadísticas (filas):"}\NormalTok{, }\FunctionTok{nrow}\NormalTok{(datos), }\StringTok{"}\SpecialCharTok{\textbackslash{}n}\StringTok{"}\NormalTok{)}
\end{Highlighting}
\end{Shaded}

\begin{verbatim}
## Unidades estadísticas (filas): 118
\end{verbatim}

\begin{Shaded}
\begin{Highlighting}[]
\FunctionTok{cat}\NormalTok{(}\StringTok{"Variables (columnas):     "}\NormalTok{, }\FunctionTok{ncol}\NormalTok{(datos), }\StringTok{"}\SpecialCharTok{\textbackslash{}n}\StringTok{"}\NormalTok{)}
\end{Highlighting}
\end{Shaded}

\begin{verbatim}
## Variables (columnas):      18
\end{verbatim}

\begin{Shaded}
\begin{Highlighting}[]
\NormalTok{utils}\SpecialCharTok{::}\FunctionTok{head}\NormalTok{(datos, }\DecValTok{10}\NormalTok{)}
\end{Highlighting}
\end{Shaded}

\begin{verbatim}
##      Cod_Div  Latitud  Longitud Region Departamento        Municipio      Fecha
## 21  50006000 3.990120 -73.76594    ORI         META          ACACÍAS 26/08/2022
## 45  50006000 3.990120 -73.76594    ORI         META          ACACÍAS 27/08/2022
## 98  85010000 5.172690 -72.54630    ORI     CASANARE          AGUAZUL 26/08/2022
## 122 85010000 5.172690 -72.54630    ORI     CASANARE          AGUAZUL 27/08/2022
## 175 81001000 7.073902 -70.74801    ORI       ARAUCA           ARAUCA 26/08/2022
## 199 81001000 7.073902 -70.74801    ORI       ARAUCA           ARAUCA 27/08/2022
## 252 81065000 7.026981 -71.42675    ORI       ARAUCA        ARAUQUITA 26/08/2022
## 276 81065000 7.026981 -71.42675    ORI       ARAUCA        ARAUQUITA 27/08/2022
## 329 50110000 4.567202 -72.96328    ORI         META BARRANCA DE UPÍA 26/08/2022
## 353 50110000 4.567202 -72.96328    ORI         META BARRANCA DE UPÍA 27/08/2022
##      Hora Temperatura Velocidad_del_Viento Direccion_del_Viento Presion
## 21  12:00        26.5                  1.2                179.8  1014.2
## 45  12:00        27.3                  1.7                113.7  1013.9
## 98  12:00        30.2                  2.1                176.7  1013.0
## 122 12:00        30.4                  1.7                200.0  1013.2
## 175 12:00        30.4                  2.0                150.0  1012.0
## 199 12:00        29.0                  0.7                187.1  1011.6
## 252 12:00        29.7                  2.1                200.7  1011.6
## 276 12:00        28.8                  0.5                356.5  1011.9
## 329 12:00        31.0                  2.0                200.0  1012.2
## 353 12:00        31.5                  0.7                210.5  1012.8
##     Punto_de_Rocio Cobertura_total_nubosa Precipitacion_mm_h
## 21            18.3                   70.0                0.0
## 45            21.0                   80.0                0.0
## 98            21.8                   49.8                0.0
## 122           24.4                   63.2                0.0
## 175           22.5                   40.0                0.0
## 199           22.3                   80.0                0.0
## 252           18.3                   50.0                0.0
## 276           24.5                   80.0                0.8
## 329           21.4                   50.0                0.0
## 353           28.7                   80.0                0.0
##     Probabilidad_de_Tormenta Humedad           Pronostico
## 21                         0    60.7 Parcialmente Nublado
## 45                         0    68.7              Nublado
## 98                         0    60.9 Parcialmente Nublado
## 122                        0    70.4 Parcialmente Nublado
## 175                        0    62.7 Parcialmente Nublado
## 199                        0    67.5              Nublado
## 252                        0    50.3 Parcialmente Nublado
## 276                        0    77.6     Nublado - Lluvia
## 329                        0    56.7 Parcialmente Nublado
## 353                        0    85.0              Nublado
\end{verbatim}

\begin{quote}
Interpretá acá qué muestra el gráfico y qué conclusión sacás. Una figura
sin lectura no es un resultado: es un adorno.
\end{quote}

\subsection{Frecuencias por clase ·
Pronostico}\label{frecuencias-por-clase-pronostico}

Panel \texttt{f1.balanceo.frecuencias} · marcado a las 07:25:51

\subsubsection{Para qué sirve}\label{para-quuxe9-sirve-1}

Ver de cuánto es el desbalance y decidir si hace falta hacer algo. La
respuesta muchas veces es que no: un 60/40 no necesita tratamiento, y
rebalancear por costumbre estropea la calibración de las probabilidades
que salen del modelo.

\subsubsection{Qué muestra}\label{quuxe9-muestra-1}

Cuántas observaciones hay por clase, antes y después de rebalancear. El
subtítulo trae la razón de desbalance: cuántas veces la clase
mayoritaria supera a la minoritaria.

Las tres técnicas disponibles no traen dependencias nuevas y no inventan
datos:

\begin{itemize}
\tightlist
\item
  \textbf{sub-muestreo} --- recorta las clases grandes. Se pierde
  información real.
\item
  \textbf{sobre-muestreo} --- repite filas de las chicas. No agrega
  información.
\item
  \textbf{bootstrap} --- remuestrea todas con reemplazo hasta un tamaño
  común.
\end{itemize}

SMOTE, que interpola entre vecinos y sí fabrica filas nuevas, está en el
catálogo como pendiente: pedirlo suma un paquete al bundle.

\subsubsection{Qué buscar}\label{quuxe9-buscar-1}

\begin{itemize}
\tightlist
\item
  \textbf{La razón de desbalance}: por debajo de 3 a 1 casi nunca hace
  falta tocar nada. Por encima de 20 a 1 el problema es real.
\item
  \textbf{Cuánto se pierde con sub-muestreo}: si la minoritaria tiene 40
  casos, la mayoritaria queda en 40 y tiraste miles de filas.
\item
  \textbf{Cuánto se repite con sobre-muestreo}: si una fila aparece 30
  veces, el modelo la va a memorizar.
\item
  \textbf{Clases con menos de 10 casos}: ninguna técnica las salva.
\end{itemize}

\subsubsection{Cuándo engaña}\label{cuuxe1ndo-engauxf1a-1}

\textbf{Balancear cambia las probabilidades a priori.} Después de
rebalancear, las probabilidades que estime el modelo no son las de la
población: están calibradas para el mundo artificial 50/50 que
construiste.

\textbf{El desbalance a veces no es el problema.} Con clases separables,
un modelo aprende bien con 1 \% de positivos. Lo que suele estar mal es
la métrica: exactitud con 99 \% de negativos se gana prediciendo siempre
``no''. Antes de remuestrear, probá cambiar de métrica o usar pesos de
clase.

\textbf{Nunca se rebalancea la partición de prueba.} Solo el
entrenamiento. Si tocás la prueba, estás midiendo el desempeño en una
población que no existe.

Estado al marcarlo:

{\def\LTcaptype{none} % do not increment counter
\begin{longtable}[]{@{}lll@{}}
\toprule\noalign{}
clase & n & proporcion \\
\midrule\noalign{}
\endhead
\bottomrule\noalign{}
\endlastfoot
Parcialmente Nublado & 63 & 0.5338983 \\
Nublado & 31 & 0.2627119 \\
Parcialmente Nublado - Lluvia Fuerte & 6 & 0.05084746 \\
Parcialmente Nublado - Llovizna & 5 & 0.04237288 \\
Nublado - Llovizna & 4 & 0.03389831 \\
Nublado - Lluvia & 4 & 0.03389831 \\
Despejado & 3 & 0.02542373 \\
Parcialmente Nublado - Tormenta Ligera & 2 & 0.01694915 \\
\end{longtable}
}

Parámetros:

{\def\LTcaptype{none} % do not increment counter
\begin{longtable}[]{@{}ll@{}}
\toprule\noalign{}
Campo & Valor \\
\midrule\noalign{}
\endhead
\bottomrule\noalign{}
\endlastfoot
clase & Pronostico \\
\end{longtable}
}

\begin{Shaded}
\begin{Highlighting}[]
\CommentTok{\# f1.balanceo.frecuencias}
\NormalTok{tabla }\OtherTok{\textless{}{-}} \FunctionTok{sort}\NormalTok{(}\FunctionTok{table}\NormalTok{(datos}\SpecialCharTok{$}\NormalTok{Pronostico), }\AttributeTok{decreasing =} \ConstantTok{TRUE}\NormalTok{)}
\FunctionTok{print}\NormalTok{(}\FunctionTok{cbind}\NormalTok{(}\AttributeTok{n =}\NormalTok{ tabla, }\AttributeTok{prop =} \FunctionTok{round}\NormalTok{(}\FunctionTok{prop.table}\NormalTok{(tabla), }\DecValTok{3}\NormalTok{)))}
\end{Highlighting}
\end{Shaded}

\begin{verbatim}
##                                         n  prop
## Parcialmente Nublado                   63 0.534
## Nublado                                31 0.263
## Parcialmente Nublado - Lluvia Fuerte    6 0.051
## Parcialmente Nublado - Llovizna         5 0.042
## Nublado - Llovizna                      4 0.034
## Nublado - Lluvia                        4 0.034
## Despejado                               3 0.025
## Parcialmente Nublado - Tormenta Ligera  2 0.017
\end{verbatim}

\begin{Shaded}
\begin{Highlighting}[]
\FunctionTok{cat}\NormalTok{(}\StringTok{"Más frecuente:  "}\NormalTok{, }\FunctionTok{names}\NormalTok{(tabla)[}\DecValTok{1}\NormalTok{], }\StringTok{" (n = "}\NormalTok{, tabla[}\DecValTok{1}\NormalTok{], }\StringTok{")}\SpecialCharTok{\textbackslash{}n}\StringTok{"}\NormalTok{,}
    \AttributeTok{sep =} \StringTok{""}\NormalTok{)}
\end{Highlighting}
\end{Shaded}

\begin{verbatim}
## Más frecuente:  Parcialmente Nublado (n = 63)
\end{verbatim}

\begin{Shaded}
\begin{Highlighting}[]
\FunctionTok{cat}\NormalTok{(}\StringTok{"Menos frecuente: "}\NormalTok{, }\FunctionTok{names}\NormalTok{(tabla)[}\FunctionTok{length}\NormalTok{(tabla)],}
    \StringTok{" (n = "}\NormalTok{, tabla[}\FunctionTok{length}\NormalTok{(tabla)], }\StringTok{")}\SpecialCharTok{\textbackslash{}n}\StringTok{"}\NormalTok{, }\AttributeTok{sep =} \StringTok{""}\NormalTok{)}
\end{Highlighting}
\end{Shaded}

\begin{verbatim}
## Menos frecuente: Parcialmente Nublado - Tormenta Ligera (n = 2)
\end{verbatim}

\begin{Shaded}
\begin{Highlighting}[]
\FunctionTok{barplot}\NormalTok{(tabla, }\AttributeTok{horiz =} \ConstantTok{TRUE}\NormalTok{, }\AttributeTok{las =} \DecValTok{1}\NormalTok{, }\AttributeTok{cex.names =} \FloatTok{0.7}\NormalTok{,}
        \AttributeTok{main =} \StringTok{"Frecuencias de Pronostico"}\NormalTok{)}
\end{Highlighting}
\end{Shaded}

\pandocbounded{\includegraphics[keepaspectratio]{taller-01_files/taller-01_files/figure-latex/unnamed-chunk-2-1.pdf}}

\begin{quote}
Interpretá acá qué muestra el gráfico y qué conclusión sacás. Una figura
sin lectura no es un resultado: es un adorno.
\end{quote}

\subsection{Diccionario de columnas}\label{diccionario-de-columnas}

Panel \texttt{f1.diccionario.tabla} · marcado a las 07:25:51

\subsubsection{Para qué sirve}\label{para-quuxe9-sirve-2}

Declarar qué es cada columna, que es lo que después habilita o bloquea
cada operación en el resto de la app. Es la card con más consecuencias
del proyecto: marcar una variable como nominal apaga su media y su
histograma a propósito, y con la razón escrita al lado.

\subsubsection{Qué muestra}\label{quuxe9-muestra-2}

\begin{Shaded}
\begin{Highlighting}[]
\NormalTok{{-}{-}{-}}
\NormalTok{config:}
\NormalTok{  layout: elk}
\NormalTok{  theme: default}
\NormalTok{{-}{-}{-}}
\NormalTok{flowchart TD}
\NormalTok{    V["CLASIFICACIÓN DE VARIABLES\textless{}br/\textgreater{}(NB1: escala × clase, ejes independientes)"]}

\NormalTok{    V {-}{-}\textgreater{} ESC["EJE 1 · ESCALA\textless{}br/\textgreater{}¿qué operaciones tienen sentido?"]}
\NormalTok{    V {-}{-}\textgreater{} CLA["EJE 2 · CLASE\textless{}br/\textgreater{}¿cuál es el recorrido?"]}

\NormalTok{    \%\% ================= EJE 1: ESCALA =================}
\NormalTok{    ESC {-}{-}\textgreater{} CUAL["CUALITATIVA\textless{}br/\textgreater{}mide atributos → categorías"]}
\NormalTok{    ESC {-}{-}\textgreater{} CUAN["CUANTITATIVA\textless{}br/\textgreater{}mide cantidades → números"]}

\NormalTok{    CUAL {-}{-}\textgreater{} NOM["NOMINAL\textless{}br/\textgreater{}solo = / ≠"]}
\NormalTok{    CUAL {-}{-}\textgreater{} ORD["ORDINAL\textless{}br/\textgreater{}= / ≠ + orden\textless{}br/\textgreater{}(distancias sin significado)"]}
\NormalTok{    CUAN {-}{-}\textgreater{} INT["INTERVALO\textless{}br/\textgreater{}+ intervalos equidistantes\textless{}br/\textgreater{}CERO RELATIVO"]}
\NormalTok{    CUAN {-}{-}\textgreater{} RAZ["RAZÓN\textless{}br/\textgreater{}+ razones válidas\textless{}br/\textgreater{}CERO ABSOLUTO"]}

\NormalTok{    NOM {-}{-}\textgreater{} NOME["Municipio, Pronóstico, Region\textless{}br/\textgreater{}(código postal: numérico pero nominal)"]}
\NormalTok{    ORD {-}{-}\textgreater{} ORDE["Estrato \{bajo, medio, alto\}\textless{}br/\textgreater{}consumo eléctrico \{alto, medio, bajo\}"]}
\NormalTok{    INT {-}{-}\textgreater{} INTE["Temperatura en °C\textless{}br/\textgreater{}0 °C no es ausencia de temperatura"]}
\NormalTok{    RAZ {-}{-}\textgreater{} RAZE["Velocidad del viento (mph)\textless{}br/\textgreater{}Presion, Latitud — 0 = ausencia"]}

\NormalTok{    \%\% ================= EJE 2: CLASE =================}
\NormalTok{    CLA {-}{-}\textgreater{} CUAL2["si es CUALITATIVA"]}
\NormalTok{    CLA {-}{-}\textgreater{} CUAN2["si es CUANTITATIVA"]}

\NormalTok{    CUAL2 {-}{-}\textgreater{} BIN["BINARIA / DICOTÓMICA\textless{}br/\textgreater{}2 categorías"]}
\NormalTok{    CUAL2 {-}{-}\textgreater{} POL["POLITÓMICA\textless{}br/\textgreater{}más de 2 categorías"]}
\NormalTok{    CUAN2 {-}{-}\textgreater{} DIS["DISCRETA\textless{}br/\textgreater{}valores contables"]}
\NormalTok{    CUAN2 {-}{-}\textgreater{} CON["CONTINUA\textless{}br/\textgreater{}cualquier valor de un intervalo real"]}

\NormalTok{    BIN {-}{-}\textgreater{} BINE["Sí / No · Presence \{Sí, No\}"]}
\NormalTok{    POL {-}{-}\textgreater{} POLE["Pronóstico (11 categorías)\textless{}br/\textgreater{}Municipio, Departamento"]}
\NormalTok{    DIS {-}{-}\textgreater{} DISE["n.° de turbinas\textless{}br/\textgreater{}n.° de hijos"]}
\NormalTok{    CON {-}{-}\textgreater{} CONE["Temperatura, Presion\textless{}br/\textgreater{}Velocidad del viento"]}

\NormalTok{    \%\% ================= ESTILOS =================}
\NormalTok{    classDef raiz fill:\#65e94f,stroke:\#222,stroke{-}width:2px,color:\#000,font{-}weight:bold;}
\NormalTok{    classDef eje fill:\#91cbe8,stroke:\#555,stroke{-}width:1px,color:\#000,font{-}weight:bold;}
\NormalTok{    classDef grupo fill:\#ffffb0,stroke:\#777,stroke{-}width:1px,color:\#000,font{-}weight:bold;}
\NormalTok{    classDef tipo fill:\#e58be8,stroke:\#777,stroke{-}width:1px,color:\#000,font{-}weight:bold;}
\NormalTok{    classDef ejemplo fill:\#8de8d8,stroke:\#777,stroke{-}width:1px,color:\#000;}

\NormalTok{    class V raiz;}
\NormalTok{    class ESC,CLA eje;}
\NormalTok{    class CUAL,CUAN,CUAL2,CUAN2 grupo;}
\NormalTok{    class NOM,ORD,INT,RAZ,BIN,POL,DIS,CON tipo;}
\NormalTok{    class NOME,ORDE,INTE,RAZE,BINE,POLE,DISE,CONE ejemplo;}
\end{Highlighting}
\end{Shaded}

Una fila por columna del dataset, con lo que la app necesita saber de
ella: etiqueta, escala de medición, clase y rol.

\begin{verbatim}
escala ∈ {nominal, ordinal, intervalo, razon}
clase  ∈ {cualitativa, discreta, continua}
rol    ∈ {respuesta, predictor, id, grupo, peso, ignorar}
\end{verbatim}

La escala no se puede deducir del tipo de dato, y ahí está todo el
punto: un código postal es numérico y es nominal; un año es entero y es
de intervalo, porque su cero es convencional. La autodetección propone;
vos corregís.

\textbf{Esta tabla gobierna el resto de la app.} Si marcás una columna
como nominal, el selector de media se deshabilita y aparece la moda, con
la razón escrita al lado. Los métodos de la fase 2 se filtran por lo
mismo.

\subsubsection{Qué buscar}\label{quuxe9-buscar-2}

\begin{itemize}
\tightlist
\item
  \textbf{Identificadores marcados como predictores}: un \texttt{id}
  correlaciona con cualquier cosa por accidente y arruina un modelo en
  silencio.
\item
  \textbf{Códigos numéricos que son categorías}: 1 = norte, 2 = sur.
  Numéricos en memoria, nominales de verdad.
\item
  \textbf{Años y fechas marcados como razón}: son de intervalo. Duplicar
  un año no significa nada.
\item
  \textbf{Columnas con muchos faltantes}: el porcentaje está en la
  tabla. Por encima de 40 \%, imputar es inventar estructura.
\item
  \textbf{Columnas constantes}: no aportan información y sobreviven
  hasta la fase 4 si nadie las marca como ignorar.
\end{itemize}

\subsubsection{Cuándo engaña}\label{cuuxe1ndo-engauxf1a-2}

\textbf{La detección automática acierta el tipo, no el significado.} Ve
números y propone razón; ve texto y propone nominal. Ordinal e intervalo
no los propone nunca, porque no se pueden inferir: los tenés que poner
vos.

\textbf{Un rol mal puesto no rompe nada hoy.} La app te deja seguir y el
error aparece dos fases después, cuando un modelo da un R²
sospechosamente perfecto porque metiste el identificador como predictor.

\textbf{Los avisos de conflicto son heurísticos.} Que una columna se
llame \texttt{year} no prueba que sea un año. El aviso pregunta; la
decisión es tuya.

\begin{Shaded}
\begin{Highlighting}[]
\CommentTok{\# f1.diccionario.tabla}
\CommentTok{\# Escala, clase y rol los declara quien analiza: R no los puede}
\CommentTok{\# deducir del tipo de dato (un código postal es numérico y nominal;}
\CommentTok{\# los grados centígrados son de intervalo, no de razón).}
\NormalTok{diccionario }\OtherTok{\textless{}{-}} \FunctionTok{data.frame}\NormalTok{(}
  \AttributeTok{columna =} \FunctionTok{c}\NormalTok{(}\StringTok{"Cod\_Div"}\NormalTok{, }\StringTok{"Latitud"}\NormalTok{, }\StringTok{"Longitud"}\NormalTok{, }\StringTok{"Region"}\NormalTok{, }\StringTok{"Departamento"}\NormalTok{,}
              \StringTok{"Municipio"}\NormalTok{, }\StringTok{"Fecha"}\NormalTok{, }\StringTok{"Hora"}\NormalTok{, }\StringTok{"Temperatura"}\NormalTok{,}
              \StringTok{"Velocidad\_del\_Viento"}\NormalTok{, }\StringTok{"Direccion\_del\_Viento"}\NormalTok{, }\StringTok{"Presion"}\NormalTok{,}
              \StringTok{"Punto\_de\_Rocio"}\NormalTok{, }\StringTok{"Cobertura\_total\_nubosa"}\NormalTok{,}
              \StringTok{"Precipitacion\_mm\_h"}\NormalTok{, }\StringTok{"Probabilidad\_de\_Tormenta"}\NormalTok{, }\StringTok{"Humedad"}\NormalTok{,}
              \StringTok{"Pronostico"}\NormalTok{),}
  \AttributeTok{escala =} \FunctionTok{c}\NormalTok{(}\StringTok{"razon"}\NormalTok{, }\StringTok{"razon"}\NormalTok{, }\StringTok{"razon"}\NormalTok{, }\StringTok{"nominal"}\NormalTok{, }\StringTok{"nominal"}\NormalTok{, }\StringTok{"nominal"}\NormalTok{,}
             \StringTok{"nominal"}\NormalTok{, }\StringTok{"nominal"}\NormalTok{, }\StringTok{"intervalo"}\NormalTok{, }\StringTok{"razon"}\NormalTok{, }\StringTok{"razon"}\NormalTok{, }\StringTok{"razon"}\NormalTok{,}
             \StringTok{"razon"}\NormalTok{, }\StringTok{"razon"}\NormalTok{, }\StringTok{"razon"}\NormalTok{, }\StringTok{"razon"}\NormalTok{, }\StringTok{"razon"}\NormalTok{, }\StringTok{"nominal"}\NormalTok{),}
  \AttributeTok{clase =} \FunctionTok{c}\NormalTok{(}\StringTok{"continua"}\NormalTok{, }\StringTok{"continua"}\NormalTok{, }\StringTok{"continua"}\NormalTok{, }\StringTok{"cualitativa"}\NormalTok{, }\StringTok{"cualitativa"}\NormalTok{,}
            \StringTok{"cualitativa"}\NormalTok{, }\StringTok{"cualitativa"}\NormalTok{, }\StringTok{"cualitativa"}\NormalTok{, }\StringTok{"continua"}\NormalTok{,}
            \StringTok{"continua"}\NormalTok{, }\StringTok{"continua"}\NormalTok{, }\StringTok{"continua"}\NormalTok{, }\StringTok{"continua"}\NormalTok{, }\StringTok{"continua"}\NormalTok{,}
            \StringTok{"continua"}\NormalTok{, }\StringTok{"discreta"}\NormalTok{, }\StringTok{"continua"}\NormalTok{, }\StringTok{"cualitativa"}\NormalTok{),}
  \AttributeTok{rol =} \FunctionTok{c}\NormalTok{(}\StringTok{"predictor"}\NormalTok{, }\StringTok{"predictor"}\NormalTok{, }\StringTok{"predictor"}\NormalTok{, }\StringTok{"predictor"}\NormalTok{, }\StringTok{"predictor"}\NormalTok{,}
          \StringTok{"predictor"}\NormalTok{, }\StringTok{"predictor"}\NormalTok{, }\StringTok{"predictor"}\NormalTok{, }\StringTok{"predictor"}\NormalTok{, }\StringTok{"predictor"}\NormalTok{,}
          \StringTok{"predictor"}\NormalTok{, }\StringTok{"predictor"}\NormalTok{, }\StringTok{"predictor"}\NormalTok{, }\StringTok{"predictor"}\NormalTok{, }\StringTok{"predictor"}\NormalTok{,}
          \StringTok{"predictor"}\NormalTok{, }\StringTok{"predictor"}\NormalTok{, }\StringTok{"predictor"}\NormalTok{)}
\NormalTok{)}
\FunctionTok{print}\NormalTok{(diccionario)}
\end{Highlighting}
\end{Shaded}

\begin{verbatim}
##                     columna    escala       clase       rol
## 1                   Cod_Div     razon    continua predictor
## 2                   Latitud     razon    continua predictor
## 3                  Longitud     razon    continua predictor
## 4                    Region   nominal cualitativa predictor
## 5              Departamento   nominal cualitativa predictor
## 6                 Municipio   nominal cualitativa predictor
## 7                     Fecha   nominal cualitativa predictor
## 8                      Hora   nominal cualitativa predictor
## 9               Temperatura intervalo    continua predictor
## 10     Velocidad_del_Viento     razon    continua predictor
## 11     Direccion_del_Viento     razon    continua predictor
## 12                  Presion     razon    continua predictor
## 13           Punto_de_Rocio     razon    continua predictor
## 14   Cobertura_total_nubosa     razon    continua predictor
## 15       Precipitacion_mm_h     razon    continua predictor
## 16 Probabilidad_de_Tormenta     razon    discreta predictor
## 17                  Humedad     razon    continua predictor
## 18               Pronostico   nominal cualitativa predictor
\end{verbatim}

\begin{Shaded}
\begin{Highlighting}[]
\CommentTok{\# Lo que R infiere solo, para comparar:}
\FunctionTok{data.frame}\NormalTok{(}\AttributeTok{columna =} \FunctionTok{names}\NormalTok{(datos),}
           \AttributeTok{clase\_en\_R =} \FunctionTok{vapply}\NormalTok{(datos, }\ControlFlowTok{function}\NormalTok{(v) }\FunctionTok{class}\NormalTok{(v)[}\DecValTok{1}\NormalTok{], }\StringTok{""}\NormalTok{),}
           \AttributeTok{distintos =} \FunctionTok{vapply}\NormalTok{(datos, }\ControlFlowTok{function}\NormalTok{(v) }\FunctionTok{length}\NormalTok{(}\FunctionTok{unique}\NormalTok{(v)), }\DecValTok{0}\DataTypeTok{L}\NormalTok{),}
           \AttributeTok{row.names =} \ConstantTok{NULL}\NormalTok{)}
\end{Highlighting}
\end{Shaded}

\begin{verbatim}
##                     columna clase_en_R distintos
## 1                   Cod_Div    integer        59
## 2                   Latitud    numeric        59
## 3                  Longitud    numeric        59
## 4                    Region  character         1
## 5              Departamento  character         8
## 6                 Municipio  character        59
## 7                     Fecha  character         2
## 8                      Hora  character         1
## 9               Temperatura    numeric        54
## 10     Velocidad_del_Viento    numeric        31
## 11     Direccion_del_Viento    numeric        91
## 12                  Presion    numeric        38
## 13           Punto_de_Rocio    numeric        79
## 14   Cobertura_total_nubosa    numeric        29
## 15       Precipitacion_mm_h    numeric        13
## 16 Probabilidad_de_Tormenta    integer         3
## 17                  Humedad    numeric       100
## 18               Pronostico  character         8
\end{verbatim}

\begin{quote}
Interpretá acá qué muestra el gráfico y qué conclusión sacás. Una figura
sin lectura no es un resultado: es un adorno.
\end{quote}

\subsection{Diagrama de caja ·
Temperatura}\label{diagrama-de-caja-temperatura}

Panel \texttt{f1.analisis.boxplot} · marcado a las 07:25:51

\subsubsection{Para qué sirve}\label{para-quuxe9-sirve-3}

Comparar posición y dispersión de un vistazo, y decidir qué variable
merece una mirada más detallada. Es el resumen que escala: veinte
boxplots en fila se leen sin esfuerzo; veinte histogramas, no.

\subsubsection{Qué muestra}\label{quuxe9-muestra-3}

El resumen de cinco números: mínimo dentro de la cerca, primer cuartil,
mediana, tercer cuartil y máximo dentro de la cerca. La caja va de Q1 a
Q3, así que encierra el 50 \% central de los datos, y la línea de
adentro es la mediana.

Los bigotes llegan hasta la observación más lejana que todavía cae
dentro de 1,5 veces el rango intercuartílico:

\[
\text{cerca inferior} = Q_1 - 1{,}5 \cdot \text{RIC}
\qquad
\text{cerca superior} = Q_3 + 1{,}5 \cdot \text{RIC}
\]

Lo que queda afuera se dibuja punto por punto. Son los \textbf{atípicos
de Tukey}, y son una convención, no un diagnóstico.

\subsubsection{Qué buscar}\label{quuxe9-buscar-3}

\begin{itemize}
\tightlist
\item
  \textbf{Posición de la mediana dentro de la caja}: si está pegada a un
  borde, la distribución es asimétrica.
\item
  \textbf{Largo relativo de los bigotes}: el más largo apunta hacia la
  cola.
\item
  \textbf{Cuántos puntos quedan afuera} y qué tan lejos. Uno muy
  separado del resto vale más atención que quince apenas pasada la
  cerca.
\item
  \textbf{Ancho de la caja}: es el RIC, la medida de dispersión que no
  se deja arrastrar por los extremos.
\end{itemize}

\subsubsection{Cuándo engaña}\label{cuuxe1ndo-engauxf1a-3}

\textbf{Esconde las modas.} Una distribución con dos picos y otra con
uno solo pueden dar cajas idénticas. El boxplot resume posiciones, no
forma. Por eso está al lado del histograma y no en su lugar.

\textbf{``Atípico'' acá significa ``pasó de 1,5·RIC'', nada más.} En una
variable con cola derecha larga ---ingresos, tiempos de espera--- es
normal que haya decenas de puntos afuera y ninguno es un error.
Marcarlos no es motivo para borrarlos.

\textbf{Con n chico los cuartiles son inestables.} Por debajo de unas 20
observaciones, agregá los puntos encima de la caja: es más honesto
mostrar el dato que su resumen.

Parámetros:

{\def\LTcaptype{none} % do not increment counter
\begin{longtable}[]{@{}ll@{}}
\toprule\noalign{}
Campo & Valor \\
\midrule\noalign{}
\endhead
\bottomrule\noalign{}
\endlastfoot
variable & Temperatura \\
\end{longtable}
}

\begin{Shaded}
\begin{Highlighting}[]
\CommentTok{\# f1.analisis.boxplot}
\FunctionTok{ggplot}\NormalTok{(datos, }\FunctionTok{aes}\NormalTok{(}\AttributeTok{x =} \StringTok{""}\NormalTok{, }\AttributeTok{y =}\NormalTok{ Temperatura)) }\SpecialCharTok{+}
  \FunctionTok{geom\_boxplot}\NormalTok{(}\AttributeTok{fill =} \StringTok{"steelblue"}\NormalTok{, }\AttributeTok{alpha =} \FloatTok{0.5}\NormalTok{) }\SpecialCharTok{+}
  \FunctionTok{coord\_flip}\NormalTok{() }\SpecialCharTok{+} \FunctionTok{labs}\NormalTok{(}\AttributeTok{x =} \ConstantTok{NULL}\NormalTok{)}
\end{Highlighting}
\end{Shaded}

\pandocbounded{\includegraphics[keepaspectratio]{taller-01_files/taller-01_files/figure-latex/unnamed-chunk-4-1.pdf}}

\begin{quote}
Interpretá acá qué muestra el gráfico y qué conclusión sacás. Una figura
sin lectura no es un resultado: es un adorno.
\end{quote}

\subsection{Diagrama de caja ·
Velocidad\_del\_Viento}\label{diagrama-de-caja-velocidad_del_viento}

Panel \texttt{f1.analisis.boxplot} · marcado a las 07:25:51

El texto de este panel está en la primera sección «Diagrama de caja».

Parámetros:

{\def\LTcaptype{none} % do not increment counter
\begin{longtable}[]{@{}ll@{}}
\toprule\noalign{}
Campo & Valor \\
\midrule\noalign{}
\endhead
\bottomrule\noalign{}
\endlastfoot
variable & Velocidad\_del\_Viento \\
\end{longtable}
}

\begin{Shaded}
\begin{Highlighting}[]
\CommentTok{\# f1.analisis.boxplot}
\FunctionTok{ggplot}\NormalTok{(datos, }\FunctionTok{aes}\NormalTok{(}\AttributeTok{x =} \StringTok{""}\NormalTok{, }\AttributeTok{y =}\NormalTok{ Velocidad\_del\_Viento)) }\SpecialCharTok{+}
  \FunctionTok{geom\_boxplot}\NormalTok{(}\AttributeTok{fill =} \StringTok{"steelblue"}\NormalTok{, }\AttributeTok{alpha =} \FloatTok{0.5}\NormalTok{) }\SpecialCharTok{+}
  \FunctionTok{coord\_flip}\NormalTok{() }\SpecialCharTok{+} \FunctionTok{labs}\NormalTok{(}\AttributeTok{x =} \ConstantTok{NULL}\NormalTok{)}
\end{Highlighting}
\end{Shaded}

\pandocbounded{\includegraphics[keepaspectratio]{taller-01_files/taller-01_files/figure-latex/unnamed-chunk-5-1.pdf}}

\begin{quote}
Interpretá acá qué muestra el gráfico y qué conclusión sacás. Una figura
sin lectura no es un resultado: es un adorno.
\end{quote}

\subsection{Resumen numérico ·
Temperatura}\label{resumen-numuxe9rico-temperatura}

Panel \texttt{f1.analisis.resumen} · marcado a las 07:25:51

\subsubsection{Para qué sirve}\label{para-quuxe9-sirve-4}

Poner en números lo que la caja muestra en forma, y comprobar que las
dos cuentan la misma historia. Es la tabla que se cita cuando hay que
escribir el centro y la dispersión de una variable sin adjetivos.

\subsubsection{Qué muestra}\label{quuxe9-muestra-4}

Los cinco números de \texttt{summary()} ---mínimo, Q1, mediana, media,
Q3, máximo--- más la desviación estándar, el rango intercuartílico y la
asimetría, con el conteo de faltantes al lado.

Qué estadísticos aparecen \textbf{lo decide la escala declarada en el
Diccionario}, no el tipo de dato en memoria:

\begin{verbatim}
nominal  -> moda y tabla de frecuencias
ordinal  -> moda, mediana y cuantiles
intervalo -> todo menos los cocientes
razon    -> todo
\end{verbatim}

La media de una nominal no sale porque no se calcula, no porque esté
escondida.

\subsubsection{Qué buscar}\label{quuxe9-buscar-4}

\begin{itemize}
\tightlist
\item
  \textbf{Media contra mediana}: si la media es menor, hay cola a la
  izquierda; si es mayor, cola a la derecha. Es la regla de bolsillo de
  la asimetría, y se confirma con el coeficiente.
\item
  \textbf{Desviación contra RIC}: cuando la desviación es mucho mayor de
  lo que el RIC sugiere, unos pocos valores extremos la están inflando.
\item
  \textbf{Dispersión relativa}: un RIC de 2,5 °C sobre una mediana de 30
  es el 8 \%; el mismo 2,5 sobre una mediana de 1,7 sería otra variable
  completamente. Comparar dispersiones en unidades distintas no
  significa nada.
\item
  \textbf{Faltantes}: cuántos son y si \texttt{summary()} los está
  descontando.
\end{itemize}

\subsubsection{Cuándo engaña}\label{cuuxe1ndo-engauxf1a-4}

\textbf{Un resumen no distingue formas.} Dos variables con los mismos
seis números pueden tener una moda o dos. Por eso esta tabla va al lado
del histograma y de la densidad, no en su lugar.

\textbf{La media de una escala de intervalo se puede calcular; el
cociente no.} 30 °C no es el doble de calor que 15 °C. El resumen no lo
avisa: lo avisa la escala que se declaró en el Diccionario.

\textbf{Con n chico los cuartiles se mueven mucho.} Por debajo de unas
20 observaciones, agregar o quitar un dato cambia Q1 y Q3 de forma
visible, y la tabla se lee con la misma desconfianza que la caja.

Estado al marcarlo:

{\def\LTcaptype{none} % do not increment counter
\begin{longtable}[]{@{}ll@{}}
\toprule\noalign{}
estadistico & mostrado \\
\midrule\noalign{}
\endhead
\bottomrule\noalign{}
\endlastfoot
n & 118 \\
faltantes & 0 \\
moda & 31 \\
mediana & 30.4 \\
q1 & 28.83 \\
q3 & 31.3 \\
rango\_inter & 2.475 \\
minimo & 22.4 \\
maximo & 32.7 \\
media & 29.75 \\
desviacion & 2.276 \\
asimetria & -1.391 \\
curtosis & 1.576 \\
\end{longtable}
}

Parámetros:

{\def\LTcaptype{none} % do not increment counter
\begin{longtable}[]{@{}ll@{}}
\toprule\noalign{}
Campo & Valor \\
\midrule\noalign{}
\endhead
\bottomrule\noalign{}
\endlastfoot
variable & Temperatura \\
escala & intervalo \\
\end{longtable}
}

\begin{Shaded}
\begin{Highlighting}[]
\CommentTok{\# f1.analisis.resumen}
\NormalTok{x }\OtherTok{\textless{}{-}}\NormalTok{ datos}\SpecialCharTok{$}\NormalTok{Temperatura}
\FunctionTok{print}\NormalTok{(}\FunctionTok{summary}\NormalTok{(x))}
\end{Highlighting}
\end{Shaded}

\begin{verbatim}
##    Min. 1st Qu.  Median    Mean 3rd Qu.    Max. 
##   22.40   28.82   30.40   29.75   31.30   32.70
\end{verbatim}

\begin{Shaded}
\begin{Highlighting}[]
\NormalTok{v }\OtherTok{\textless{}{-}}\NormalTok{ x[}\SpecialCharTok{!}\FunctionTok{is.na}\NormalTok{(x)]}
\NormalTok{centrado }\OtherTok{\textless{}{-}}\NormalTok{ v }\SpecialCharTok{{-}} \FunctionTok{mean}\NormalTok{(v)}
\NormalTok{raiz }\OtherTok{\textless{}{-}} \FunctionTok{sqrt}\NormalTok{(}\FunctionTok{mean}\NormalTok{(centrado}\SpecialCharTok{\^{}}\DecValTok{2}\NormalTok{))}
\NormalTok{asimetria }\OtherTok{\textless{}{-}} \FunctionTok{mean}\NormalTok{(centrado}\SpecialCharTok{\^{}}\DecValTok{3}\NormalTok{) }\SpecialCharTok{/}\NormalTok{ raiz}\SpecialCharTok{\^{}}\DecValTok{3}   \CommentTok{\# g1: hacia dónde va la cola}
\NormalTok{curtosis }\OtherTok{\textless{}{-}} \FunctionTok{mean}\NormalTok{(centrado}\SpecialCharTok{\^{}}\DecValTok{4}\NormalTok{) }\SpecialCharTok{/}\NormalTok{ raiz}\SpecialCharTok{\^{}}\DecValTok{4} \SpecialCharTok{{-}} \DecValTok{3}  \CommentTok{\# g2: exceso sobre la normal}
\FunctionTok{cat}\NormalTok{(}\StringTok{"Desviación:"}\NormalTok{, }\FunctionTok{round}\NormalTok{(}\FunctionTok{sd}\NormalTok{(v), }\DecValTok{4}\NormalTok{), }\StringTok{"· RIC:"}\NormalTok{, }\FunctionTok{round}\NormalTok{(}\FunctionTok{IQR}\NormalTok{(v), }\DecValTok{4}\NormalTok{), }\StringTok{"}\SpecialCharTok{\textbackslash{}n}\StringTok{"}\NormalTok{)}
\end{Highlighting}
\end{Shaded}

\begin{verbatim}
## Desviación: 2.2756 · RIC: 2.475
\end{verbatim}

\begin{Shaded}
\begin{Highlighting}[]
\FunctionTok{cat}\NormalTok{(}\StringTok{"Asimetría:"}\NormalTok{, }\FunctionTok{round}\NormalTok{(asimetria, }\DecValTok{3}\NormalTok{),}
    \StringTok{"· Curtosis:"}\NormalTok{, }\FunctionTok{round}\NormalTok{(curtosis, }\DecValTok{3}\NormalTok{), }\StringTok{"}\SpecialCharTok{\textbackslash{}n}\StringTok{"}\NormalTok{)}
\end{Highlighting}
\end{Shaded}

\begin{verbatim}
## Asimetría: -1.391 · Curtosis: 1.576
\end{verbatim}

\begin{Shaded}
\begin{Highlighting}[]
\FunctionTok{cat}\NormalTok{(}\StringTok{"n:"}\NormalTok{, }\FunctionTok{length}\NormalTok{(v), }\StringTok{"· faltantes:"}\NormalTok{, }\FunctionTok{sum}\NormalTok{(}\FunctionTok{is.na}\NormalTok{(x)), }\StringTok{"}\SpecialCharTok{\textbackslash{}n}\StringTok{"}\NormalTok{)}
\end{Highlighting}
\end{Shaded}

\begin{verbatim}
## n: 118 · faltantes: 0
\end{verbatim}

\begin{quote}
Interpretá acá qué muestra el gráfico y qué conclusión sacás. Una figura
sin lectura no es un resultado: es un adorno.
\end{quote}

\subsection{Resumen numérico ·
Velocidad\_del\_Viento}\label{resumen-numuxe9rico-velocidad_del_viento}

Panel \texttt{f1.analisis.resumen} · marcado a las 07:25:51

El texto de este panel está en la primera sección «Resumen numérico».

Estado al marcarlo:

{\def\LTcaptype{none} % do not increment counter
\begin{longtable}[]{@{}ll@{}}
\toprule\noalign{}
estadistico & mostrado \\
\midrule\noalign{}
\endhead
\bottomrule\noalign{}
\endlastfoot
n & 118 \\
faltantes & 0 \\
moda & 0.9 \\
mediana & 1.7 \\
q1 & 1 \\
q3 & 2.175 \\
rango\_inter & 1.175 \\
minimo & 0.4 \\
maximo & 3.7 \\
media & 1.719 \\
desviacion & 0.7892 \\
asimetria & 0.4854 \\
curtosis & -0.3234 \\
coef\_variacion & 0.4592 \\
\end{longtable}
}

Parámetros:

{\def\LTcaptype{none} % do not increment counter
\begin{longtable}[]{@{}ll@{}}
\toprule\noalign{}
Campo & Valor \\
\midrule\noalign{}
\endhead
\bottomrule\noalign{}
\endlastfoot
variable & Velocidad\_del\_Viento \\
escala & razon \\
\end{longtable}
}

\begin{Shaded}
\begin{Highlighting}[]
\CommentTok{\# f1.analisis.resumen}
\NormalTok{x }\OtherTok{\textless{}{-}}\NormalTok{ datos}\SpecialCharTok{$}\NormalTok{Velocidad\_del\_Viento}
\FunctionTok{print}\NormalTok{(}\FunctionTok{summary}\NormalTok{(x))}
\end{Highlighting}
\end{Shaded}

\begin{verbatim}
##    Min. 1st Qu.  Median    Mean 3rd Qu.    Max. 
##   0.400   1.000   1.700   1.719   2.175   3.700
\end{verbatim}

\begin{Shaded}
\begin{Highlighting}[]
\NormalTok{v }\OtherTok{\textless{}{-}}\NormalTok{ x[}\SpecialCharTok{!}\FunctionTok{is.na}\NormalTok{(x)]}
\NormalTok{centrado }\OtherTok{\textless{}{-}}\NormalTok{ v }\SpecialCharTok{{-}} \FunctionTok{mean}\NormalTok{(v)}
\NormalTok{raiz }\OtherTok{\textless{}{-}} \FunctionTok{sqrt}\NormalTok{(}\FunctionTok{mean}\NormalTok{(centrado}\SpecialCharTok{\^{}}\DecValTok{2}\NormalTok{))}
\NormalTok{asimetria }\OtherTok{\textless{}{-}} \FunctionTok{mean}\NormalTok{(centrado}\SpecialCharTok{\^{}}\DecValTok{3}\NormalTok{) }\SpecialCharTok{/}\NormalTok{ raiz}\SpecialCharTok{\^{}}\DecValTok{3}   \CommentTok{\# g1: hacia dónde va la cola}
\NormalTok{curtosis }\OtherTok{\textless{}{-}} \FunctionTok{mean}\NormalTok{(centrado}\SpecialCharTok{\^{}}\DecValTok{4}\NormalTok{) }\SpecialCharTok{/}\NormalTok{ raiz}\SpecialCharTok{\^{}}\DecValTok{4} \SpecialCharTok{{-}} \DecValTok{3}  \CommentTok{\# g2: exceso sobre la normal}
\FunctionTok{cat}\NormalTok{(}\StringTok{"Desviación:"}\NormalTok{, }\FunctionTok{round}\NormalTok{(}\FunctionTok{sd}\NormalTok{(v), }\DecValTok{4}\NormalTok{), }\StringTok{"· RIC:"}\NormalTok{, }\FunctionTok{round}\NormalTok{(}\FunctionTok{IQR}\NormalTok{(v), }\DecValTok{4}\NormalTok{), }\StringTok{"}\SpecialCharTok{\textbackslash{}n}\StringTok{"}\NormalTok{)}
\end{Highlighting}
\end{Shaded}

\begin{verbatim}
## Desviación: 0.7892 · RIC: 1.175
\end{verbatim}

\begin{Shaded}
\begin{Highlighting}[]
\FunctionTok{cat}\NormalTok{(}\StringTok{"Asimetría:"}\NormalTok{, }\FunctionTok{round}\NormalTok{(asimetria, }\DecValTok{3}\NormalTok{),}
    \StringTok{"· Curtosis:"}\NormalTok{, }\FunctionTok{round}\NormalTok{(curtosis, }\DecValTok{3}\NormalTok{), }\StringTok{"}\SpecialCharTok{\textbackslash{}n}\StringTok{"}\NormalTok{)}
\end{Highlighting}
\end{Shaded}

\begin{verbatim}
## Asimetría: 0.485 · Curtosis: -0.323
\end{verbatim}

\begin{Shaded}
\begin{Highlighting}[]
\FunctionTok{cat}\NormalTok{(}\StringTok{"n:"}\NormalTok{, }\FunctionTok{length}\NormalTok{(v), }\StringTok{"· faltantes:"}\NormalTok{, }\FunctionTok{sum}\NormalTok{(}\FunctionTok{is.na}\NormalTok{(x)), }\StringTok{"}\SpecialCharTok{\textbackslash{}n}\StringTok{"}\NormalTok{)}
\end{Highlighting}
\end{Shaded}

\begin{verbatim}
## n: 118 · faltantes: 0
\end{verbatim}

\begin{quote}
Interpretá acá qué muestra el gráfico y qué conclusión sacás. Una figura
sin lectura no es un resultado: es un adorno.
\end{quote}

\subsection{Atípicos detectados ·
Presion}\label{atuxedpicos-detectados-presion}

Panel \texttt{f1.calidad.atipicos} · marcado a las 07:25:51

\subsubsection{Para qué sirve}\label{para-quuxe9-sirve-5}

Decidir qué hacer con los valores extremos: corregirlos, dejarlos, o
cambiar a un método robusto que no dependa de ellos. Lo que no se decide
acá se decide solo más adelante, y casi siempre mal.

\subsubsection{Qué muestra}\label{quuxe9-muestra-5}

Cada observación contra su posición en el archivo, con los atípicos
marcados según el criterio elegido. Los tres criterios responden
preguntas distintas:

\[
\begin{aligned}
\text{IQR} &: && x \notin \left[\, Q_1 - 1.5 \cdot \text{RIC} \;,\; Q_3 + 1.5 \cdot \text{RIC} \,\right] && \text{robusto, univariado} \\[4pt]
z &: && \frac{\lvert x - \bar{x} \rvert}{s} > 3 && \text{no robusto, univariado} \\[4pt]
\text{Mahalanobis} &: && d^2 > \chi^2_p(\alpha) && \text{multivariado}
\end{aligned}
\]

Los dos primeros miran una columna; el tercero mira todas a la vez y
encuentra filas que ninguna columna por separado delata.

\subsubsection{Qué buscar}\label{quuxe9-buscar-5}

\begin{itemize}
\tightlist
\item
  \textbf{Puntos aislados muy lejos del resto}: candidatos a error de
  captura.
\item
  \textbf{Grupos de atípicos juntos en el eje horizontal}: no son ruido,
  es un tramo del archivo distinto ---otro año, otro instrumento, otro
  criterio de medición.
\item
  \textbf{Muchos atípicos de un solo lado}: no hay error, hay asimetría.
  La transformación log suele hacerlos desaparecer sin borrar nada.
\item
  \textbf{La diferencia entre criterios}: si z no marca nada y el RIC
  marca veinte, es porque la media y la desviación ya fueron arrastradas
  por los extremos.
\end{itemize}

\subsubsection{Cuándo engaña}\label{cuuxe1ndo-engauxf1a-5}

\textbf{Atípico no es sinónimo de error.} El criterio es geométrico y no
sabe nada del fenómeno. Borrar por costumbre es la peor decisión
posible: en muchos estudios el atípico es el caso que motivó la
investigación.

\textbf{z se enmascara a sí mismo.} La media y la desviación estándar se
calculan con los atípicos adentro, así que un valor extremo infla
\texttt{s} y termina pareciendo razonable. Con un solo extremo muy
grande, ningún punto supera 3 desvíos.

\textbf{El 1,5 de Tukey es una convención, no un umbral óptimo.} Nada se
rompe si lo movés; solo cambia cuántos puntos se pintan.

\textbf{Mahalanobis exige covarianza invertible y datos completos.} Con
colinealidad exacta o con filas incompletas devuelve vacío, y eso no
significa que no haya atípicos.

Parámetros:

{\def\LTcaptype{none} % do not increment counter
\begin{longtable}[]{@{}ll@{}}
\toprule\noalign{}
Campo & Valor \\
\midrule\noalign{}
\endhead
\bottomrule\noalign{}
\endlastfoot
columna & Presion \\
\end{longtable}
}

\begin{Shaded}
\begin{Highlighting}[]
\CommentTok{\# f1.calidad.atipicos}
\NormalTok{x }\OtherTok{\textless{}{-}}\NormalTok{ datos}\SpecialCharTok{$}\NormalTok{Presion}
\CommentTok{\# Ojo: boxplot() corta por bisagras y el panel del lab por cuartiles}
\CommentTok{\# (tipo 7, que es con lo que ggplot2 dibuja la caja). Las cercas pueden}
\CommentTok{\# diferir en décimas; el conteo de atípicos casi nunca cambia.}
\NormalTok{caja }\OtherTok{\textless{}{-}} \FunctionTok{boxplot}\NormalTok{(x, }\AttributeTok{plot =} \ConstantTok{FALSE}\NormalTok{)}
\NormalTok{atipicos }\OtherTok{\textless{}{-}}\NormalTok{ caja}\SpecialCharTok{$}\NormalTok{out}
\NormalTok{ric }\OtherTok{\textless{}{-}} \FunctionTok{diff}\NormalTok{(caja}\SpecialCharTok{$}\NormalTok{stats[}\FunctionTok{c}\NormalTok{(}\DecValTok{2}\NormalTok{, }\DecValTok{4}\NormalTok{)])}
\NormalTok{cercas }\OtherTok{\textless{}{-}} \FunctionTok{c}\NormalTok{(caja}\SpecialCharTok{$}\NormalTok{stats[}\DecValTok{2}\NormalTok{] }\SpecialCharTok{{-}} \FloatTok{1.5} \SpecialCharTok{*}\NormalTok{ ric, caja}\SpecialCharTok{$}\NormalTok{stats[}\DecValTok{4}\NormalTok{] }\SpecialCharTok{+} \FloatTok{1.5} \SpecialCharTok{*}\NormalTok{ ric)}
\FunctionTok{cat}\NormalTok{(}\StringTok{"Atípicos por el criterio de Tukey:"}\NormalTok{, }\FunctionTok{length}\NormalTok{(atipicos),}
    \StringTok{"de"}\NormalTok{, }\FunctionTok{length}\NormalTok{(x), }\StringTok{"}\SpecialCharTok{\textbackslash{}n}\StringTok{"}\NormalTok{)}
\end{Highlighting}
\end{Shaded}

\begin{verbatim}
## Atípicos por el criterio de Tukey: 1 de 118
\end{verbatim}

\begin{Shaded}
\begin{Highlighting}[]
\FunctionTok{cat}\NormalTok{(}\StringTok{"Cercas:"}\NormalTok{, }\FunctionTok{round}\NormalTok{(cercas, }\DecValTok{2}\NormalTok{), }\StringTok{"}\SpecialCharTok{\textbackslash{}n}\StringTok{"}\NormalTok{)}
\end{Highlighting}
\end{Shaded}

\begin{verbatim}
## Cercas: 1009.65 1015.65
\end{verbatim}

\begin{Shaded}
\begin{Highlighting}[]
\FunctionTok{print}\NormalTok{(atipicos)}
\end{Highlighting}
\end{Shaded}

\begin{verbatim}
## [1] 1015.7
\end{verbatim}

\begin{Shaded}
\begin{Highlighting}[]
\FunctionTok{boxplot}\NormalTok{(x, }\AttributeTok{horizontal =} \ConstantTok{TRUE}\NormalTok{, }\AttributeTok{main =} \StringTok{"Presion"}\NormalTok{)}
\end{Highlighting}
\end{Shaded}

\pandocbounded{\includegraphics[keepaspectratio]{taller-01_files/taller-01_files/figure-latex/unnamed-chunk-8-1.pdf}}

\begin{quote}
Interpretá acá qué muestra el gráfico y qué conclusión sacás. Una figura
sin lectura no es un resultado: es un adorno.
\end{quote}

\subsection{Atípicos detectados ·
Punto\_de\_Rocio}\label{atuxedpicos-detectados-punto_de_rocio}

Panel \texttt{f1.calidad.atipicos} · marcado a las 07:25:51

El texto de este panel está en la primera sección «Atípicos detectados».

Parámetros:

{\def\LTcaptype{none} % do not increment counter
\begin{longtable}[]{@{}ll@{}}
\toprule\noalign{}
Campo & Valor \\
\midrule\noalign{}
\endhead
\bottomrule\noalign{}
\endlastfoot
columna & Punto\_de\_Rocio \\
\end{longtable}
}

\begin{Shaded}
\begin{Highlighting}[]
\CommentTok{\# f1.calidad.atipicos}
\NormalTok{x }\OtherTok{\textless{}{-}}\NormalTok{ datos}\SpecialCharTok{$}\NormalTok{Punto\_de\_Rocio}
\CommentTok{\# Ojo: boxplot() corta por bisagras y el panel del lab por cuartiles}
\CommentTok{\# (tipo 7, que es con lo que ggplot2 dibuja la caja). Las cercas pueden}
\CommentTok{\# diferir en décimas; el conteo de atípicos casi nunca cambia.}
\NormalTok{caja }\OtherTok{\textless{}{-}} \FunctionTok{boxplot}\NormalTok{(x, }\AttributeTok{plot =} \ConstantTok{FALSE}\NormalTok{)}
\NormalTok{atipicos }\OtherTok{\textless{}{-}}\NormalTok{ caja}\SpecialCharTok{$}\NormalTok{out}
\NormalTok{ric }\OtherTok{\textless{}{-}} \FunctionTok{diff}\NormalTok{(caja}\SpecialCharTok{$}\NormalTok{stats[}\FunctionTok{c}\NormalTok{(}\DecValTok{2}\NormalTok{, }\DecValTok{4}\NormalTok{)])}
\NormalTok{cercas }\OtherTok{\textless{}{-}} \FunctionTok{c}\NormalTok{(caja}\SpecialCharTok{$}\NormalTok{stats[}\DecValTok{2}\NormalTok{] }\SpecialCharTok{{-}} \FloatTok{1.5} \SpecialCharTok{*}\NormalTok{ ric, caja}\SpecialCharTok{$}\NormalTok{stats[}\DecValTok{4}\NormalTok{] }\SpecialCharTok{+} \FloatTok{1.5} \SpecialCharTok{*}\NormalTok{ ric)}
\FunctionTok{cat}\NormalTok{(}\StringTok{"Atípicos por el criterio de Tukey:"}\NormalTok{, }\FunctionTok{length}\NormalTok{(atipicos),}
    \StringTok{"de"}\NormalTok{, }\FunctionTok{length}\NormalTok{(x), }\StringTok{"}\SpecialCharTok{\textbackslash{}n}\StringTok{"}\NormalTok{)}
\end{Highlighting}
\end{Shaded}

\begin{verbatim}
## Atípicos por el criterio de Tukey: 0 de 118
\end{verbatim}

\begin{Shaded}
\begin{Highlighting}[]
\FunctionTok{cat}\NormalTok{(}\StringTok{"Cercas:"}\NormalTok{, }\FunctionTok{round}\NormalTok{(cercas, }\DecValTok{2}\NormalTok{), }\StringTok{"}\SpecialCharTok{\textbackslash{}n}\StringTok{"}\NormalTok{)}
\end{Highlighting}
\end{Shaded}

\begin{verbatim}
## Cercas: 13.35 30.55
\end{verbatim}

\begin{Shaded}
\begin{Highlighting}[]
\FunctionTok{print}\NormalTok{(atipicos)}
\end{Highlighting}
\end{Shaded}

\begin{verbatim}
## numeric(0)
\end{verbatim}

\begin{Shaded}
\begin{Highlighting}[]
\FunctionTok{boxplot}\NormalTok{(x, }\AttributeTok{horizontal =} \ConstantTok{TRUE}\NormalTok{, }\AttributeTok{main =} \StringTok{"Punto\_de\_Rocio"}\NormalTok{)}
\end{Highlighting}
\end{Shaded}

\pandocbounded{\includegraphics[keepaspectratio]{taller-01_files/taller-01_files/figure-latex/unnamed-chunk-9-1.pdf}}

\begin{quote}
Interpretá acá qué muestra el gráfico y qué conclusión sacás. Una figura
sin lectura no es un resultado: es un adorno.
\end{quote}

\subsection{Diagrama de dispersión · Temperatura
vs.~Velocidad\_del\_Viento}\label{diagrama-de-dispersiuxf3n-temperatura-vs.-velocidad_del_viento}

Panel \texttt{f1.analisis.dispersion} · marcado a las 07:25:51

\subsubsection{Para qué sirve}\label{para-quuxe9-sirve-6}

Decidir si dos variables tienen relación, y de qué tipo, antes de
resumirla en un coeficiente. Todo lo que venga después ---correlación,
regresión--- da por hecha una forma; acá se comprueba que esa forma
existe.

\subsubsection{Qué muestra}\label{quuxe9-muestra-6}

Cada observación como un punto en el plano de dos variables. Es el
gráfico más directo que existe y el que sostiene toda la idea de
regresión: si hay relación, se ve acá antes de ajustar nada.

El subtítulo trae dos números que conviene leer juntos:

\begin{verbatim}
r   = correlación de Pearson  → mide relación LINEAL
rho = correlación de Spearman → mide relación MONÓTONA
\end{verbatim}

Cuando \texttt{\textbar{}rho\textbar{}} es bastante mayor que
\texttt{\textbar{}r\textbar{}}, hay relación y no es una recta.

\subsubsection{Qué buscar}\label{quuxe9-buscar-6}

\begin{itemize}
\tightlist
\item
  \textbf{Forma}: recta, curva, abanico, nube sin estructura.
\item
  \textbf{Dirección y fuerza}: cuánto se aprieta la nube alrededor de su
  tendencia.
\item
  \textbf{Heterocedasticidad}: si la dispersión vertical crece con x, el
  abanico está avisando que la varianza no es constante.
\item
  \textbf{Grupos}: dos nubes separadas piden colorear por una tercera
  variable.
\item
  \textbf{Puntos con influencia}: los que están lejos en x mueven una
  recta ajustada mucho más que los que están lejos en y.
\end{itemize}

\subsubsection{Cuándo engaña}\label{cuuxe1ndo-engauxf1a-6}

\textbf{El sobreploteo miente por saturación.} Con muchos puntos, la
mancha negra oculta dónde está la masa: cualquier zona llena se ve igual
de negra que una diez veces más densa. Las tres curas están en el panel:
bajar la transparencia, agregar jitter cuando los valores están
redondeados, o pasar a conteo por celda, que ya no dibuja puntos sino
cuánta gente hay en cada casilla.

\textbf{r = 0 no significa ``sin relación''.} Significa ``sin relación
lineal''. Una parábola perfecta da r cercano a cero. Mirá la nube, no el
número.

\textbf{Correlación no es causa, y encima puede ser de un tercero.} Dos
variables pueden moverse juntas porque una tercera las mueve a las dos.
Colorear por esa tercera suele deshacer el espejismo en un segundo.

\textbf{El muestreo cambia el dibujo, no la cuenta.} Si el badge de
muestreo está activo, los puntos son una muestra con semilla; \texttt{r}
y \texttt{rho} se calcularon sobre el total.

Parámetros:

{\def\LTcaptype{none} % do not increment counter
\begin{longtable}[]{@{}ll@{}}
\toprule\noalign{}
Campo & Valor \\
\midrule\noalign{}
\endhead
\bottomrule\noalign{}
\endlastfoot
x & Temperatura \\
y & Velocidad\_del\_Viento \\
marginales & TRUE \\
\end{longtable}
}

\begin{Shaded}
\begin{Highlighting}[]
\CommentTok{\# f1.analisis.dispersion}
\NormalTok{grafico }\OtherTok{\textless{}{-}} \FunctionTok{ggplot}\NormalTok{(datos, }\FunctionTok{aes}\NormalTok{(}\AttributeTok{x =}\NormalTok{ Temperatura, }\AttributeTok{y =}\NormalTok{ Velocidad\_del\_Viento)) }\SpecialCharTok{+}
  \FunctionTok{geom\_point}\NormalTok{(}\AttributeTok{alpha =} \FloatTok{0.6}\NormalTok{) }\SpecialCharTok{+}
  \FunctionTok{labs}\NormalTok{(}\AttributeTok{x =} \StringTok{"Temperatura"}\NormalTok{, }\AttributeTok{y =} \StringTok{"Velocidad\_del\_Viento"}\NormalTok{)}
\NormalTok{ggExtra}\SpecialCharTok{::}\FunctionTok{ggMarginal}\NormalTok{(grafico, }\AttributeTok{type =} \StringTok{"histogram"}\NormalTok{, }\AttributeTok{fill =} \StringTok{"grey70"}\NormalTok{,}
                    \AttributeTok{color =} \StringTok{"white"}\NormalTok{)}
\end{Highlighting}
\end{Shaded}

\pandocbounded{\includegraphics[keepaspectratio]{taller-01_files/taller-01_files/figure-latex/unnamed-chunk-10-1.pdf}}

\begin{Shaded}
\begin{Highlighting}[]
\FunctionTok{cat}\NormalTok{(}\StringTok{"Covarianza:"}\NormalTok{, }\FunctionTok{cov}\NormalTok{(datos}\SpecialCharTok{$}\NormalTok{Temperatura, datos}\SpecialCharTok{$}\NormalTok{Velocidad\_del\_Viento, }\AttributeTok{use =} \StringTok{"complete.obs"}\NormalTok{), }\StringTok{"}\SpecialCharTok{\textbackslash{}n}\StringTok{"}\NormalTok{)}
\end{Highlighting}
\end{Shaded}

\begin{verbatim}
## Covarianza: 0.04718528
\end{verbatim}

\begin{Shaded}
\begin{Highlighting}[]
\FunctionTok{cat}\NormalTok{(}\StringTok{"Pearson:   "}\NormalTok{, }\FunctionTok{cor}\NormalTok{(datos}\SpecialCharTok{$}\NormalTok{Temperatura, datos}\SpecialCharTok{$}\NormalTok{Velocidad\_del\_Viento, }\AttributeTok{use =} \StringTok{"complete.obs"}\NormalTok{), }\StringTok{"}\SpecialCharTok{\textbackslash{}n}\StringTok{"}\NormalTok{)}
\end{Highlighting}
\end{Shaded}

\begin{verbatim}
## Pearson:    0.02627245
\end{verbatim}

\begin{Shaded}
\begin{Highlighting}[]
\FunctionTok{cat}\NormalTok{(}\StringTok{"Spearman:  "}\NormalTok{, }\FunctionTok{cor}\NormalTok{(datos}\SpecialCharTok{$}\NormalTok{Temperatura, datos}\SpecialCharTok{$}\NormalTok{Velocidad\_del\_Viento, }\AttributeTok{use =} \StringTok{"complete.obs"}\NormalTok{,}
    \AttributeTok{method =} \StringTok{"spearman"}\NormalTok{), }\StringTok{"}\SpecialCharTok{\textbackslash{}n}\StringTok{"}\NormalTok{)}
\end{Highlighting}
\end{Shaded}

\begin{verbatim}
## Spearman:   -0.01089905
\end{verbatim}

\begin{Shaded}
\begin{Highlighting}[]
\FunctionTok{print}\NormalTok{(}\FunctionTok{cor.test}\NormalTok{(datos}\SpecialCharTok{$}\NormalTok{Temperatura, datos}\SpecialCharTok{$}\NormalTok{Velocidad\_del\_Viento))}
\end{Highlighting}
\end{Shaded}

\begin{verbatim}
## 
##  Pearson's product-moment correlation
## 
## data:  datos$Temperatura and datos$Velocidad_del_Viento
## t = 0.28306, df = 116, p-value = 0.7776
## alternative hypothesis: true correlation is not equal to 0
## 95 percent confidence interval:
##  -0.1552241  0.2060533
## sample estimates:
##        cor 
## 0.02627245
\end{verbatim}

\begin{quote}
Interpretá acá qué muestra el gráfico y qué conclusión sacás. Una figura
sin lectura no es un resultado: es un adorno.
\end{quote}

\subsection{Histograma · Temperatura}\label{histograma-temperatura}

Panel \texttt{f1.analisis.histograma} · marcado a las 07:25:51

\subsubsection{Para qué sirve}\label{para-quuxe9-sirve-7}

Formarse la primera idea de una variable antes de decidir cualquier otra
cosa sobre ella: si transformarla, si partirla por grupos, si desconfiar
de cómo se midió. Es la card que se abre primero y la que más veces
termina cambiando el plan.

\subsubsection{Qué muestra}\label{quuxe9-muestra-7}

Cuántas observaciones caen en cada intervalo de la variable. La altura
de cada barra es una frecuencia; el ancho, un tramo del recorrido.

Es el primer retrato de una variable continua y el antecedente directo
de la idea de \textbf{densidad}: si en vez de contar dividís por el
total y por el ancho del intervalo, el área total pasa a valer 1 y ya
estás mirando una densidad empírica.

\subsubsection{Qué buscar}\label{quuxe9-buscar-7}

\begin{itemize}
\tightlist
\item
  \textbf{Centro}: ¿dónde está el grueso de los datos? ¿Coincide con la
  media?
\item
  \textbf{Dispersión}: ¿ancho o angosto respecto al recorrido posible?
\item
  \textbf{Forma}: ¿simétrico o con cola a un lado? Una cola derecha
  larga es la firma de variables que no pueden ser negativas (ingresos,
  tiempos, conteos), y suele pedir una transformación logarítmica.
\item
  \textbf{Modas}: ¿un pico o varios? Dos picos casi siempre significan
  que hay dos poblaciones mezcladas, y ese es un hallazgo, no un
  defecto.
\item
  \textbf{Huecos y paredes}: un vacío en medio, o un corte abrupto en un
  valor redondo, suelen delatar cómo se recogió el dato, no cómo se
  comporta.
\end{itemize}

\subsubsection{Cuándo engaña}\label{cuuxe1ndo-engauxf1a-7}

\textbf{El número de clases cambia la historia.} Es la trampa central de
este gráfico y por eso el control está a la vista. Con pocas barras todo
parece una campana suave; con muchas, todo parece ruido dentado. Ninguna
de las dos versiones es ``la verdadera'': son la misma tabla de datos
leída con distinta resolución.

Movelo. Si una moda secundaria aparece y desaparece según el número de
clases, no confíes en ella todavía --- comparala con la densidad kernel,
que suaviza sin depender de dónde caen los bordes.

\textbf{El origen de los intervalos también importa.} Dos histogramas
con el mismo ancho de clase pero desplazados medio intervalo pueden
verse distintos. La densidad kernel no tiene ese problema, porque no
tiene bordes.

\textbf{Con pocos datos no dice casi nada.} Por debajo de unas 30
observaciones, la forma es principalmente azar. Ahí conviene un diagrama
de puntos o el de tallo y hojas, que no esconden ninguna observación.

Parámetros:

{\def\LTcaptype{none} % do not increment counter
\begin{longtable}[]{@{}ll@{}}
\toprule\noalign{}
Campo & Valor \\
\midrule\noalign{}
\endhead
\bottomrule\noalign{}
\endlastfoot
variable & Temperatura \\
densidad & TRUE \\
\end{longtable}
}

\begin{Shaded}
\begin{Highlighting}[]
\CommentTok{\# f1.analisis.histograma}
\FunctionTok{ggplot}\NormalTok{(datos, }\FunctionTok{aes}\NormalTok{(}\AttributeTok{x =}\NormalTok{ Temperatura)) }\SpecialCharTok{+}
  \FunctionTok{geom\_histogram}\NormalTok{(}\FunctionTok{aes}\NormalTok{(}\AttributeTok{y =} \FunctionTok{after\_stat}\NormalTok{(density)),}
                 \AttributeTok{bins =} \DecValTok{30}\NormalTok{, }\AttributeTok{fill =} \StringTok{"steelblue"}\NormalTok{) }\SpecialCharTok{+}
  \FunctionTok{geom\_density}\NormalTok{(}\AttributeTok{alpha =} \FloatTok{0.4}\NormalTok{) }\SpecialCharTok{+}
  \FunctionTok{labs}\NormalTok{(}\AttributeTok{x =} \StringTok{"Temperatura"}\NormalTok{, }\AttributeTok{y =} \StringTok{"densidad"}\NormalTok{)}
\end{Highlighting}
\end{Shaded}

\pandocbounded{\includegraphics[keepaspectratio]{taller-01_files/taller-01_files/figure-latex/unnamed-chunk-11-1.pdf}}

\begin{quote}
Interpretá acá qué muestra el gráfico y qué conclusión sacás. Una figura
sin lectura no es un resultado: es un adorno.
\end{quote}

\subsection{Q-Q normal de la variable ·
Temperatura}\label{q-q-normal-de-la-variable-temperatura}

Panel \texttt{f1.analisis.qq\_normal\_datos} · marcado a las 07:25:51

\subsubsection{Para qué sirve}\label{para-quuxe9-sirve-8}

Comprobar la normalidad de una variable con más resolución que un
histograma, sobre todo en las colas. Es donde se distingue si el
problema son unos pocos valores extremos o la forma completa, y esa
diferencia decide si alcanza con recortar o hay que cambiar de método.

\subsubsection{Qué muestra}\label{quuxe9-muestra-8}

Los cuantiles observados de la variable contra los que tendría si fuera
normal. Cada punto es una observación: en el eje horizontal, el valor
que le tocaría bajo una normal; en el vertical, el que de verdad tiene.

Si la variable es normal, los puntos caen sobre la recta de referencia,
que pasa por el primer y el tercer cuartil. Los desvíos se leen por su
forma, no por su tamaño.

\subsubsection{Qué buscar}\label{quuxe9-buscar-8}

\begin{itemize}
\tightlist
\item
  \textbf{S acostada}: asimetría. Si los puntos se despegan hacia arriba
  en la derecha, hay cola derecha larga.
\item
  \textbf{Extremos que se van hacia afuera en los dos lados}: colas
  pesadas, más masa lejos del centro de la que una normal admite.
\item
  \textbf{Extremos que se aplanan}: colas livianas, distribución más
  corta que la normal.
\item
  \textbf{Escalones horizontales}: la variable está discretizada o
  redondeada.
\item
  \textbf{Uno o dos puntos muy lejos de la recta al final}: atípicos, no
  forma.
\end{itemize}

\subsubsection{Cuándo engaña}\label{cuuxe1ndo-engauxf1a-8}

\textbf{La prueba y el gráfico se contradicen a propósito, y el gráfico
tiene razón más veces.} Con n grande, Shapiro-Wilk rechaza la normalidad
por desvíos minúsculos que no afectan a ningún método; con n chico, no
rechaza nada porque no tiene potencia. El subtítulo trae la prueba y el
p, pero la decisión se toma mirando la nube.

\textbf{Casi ningún método pide que los DATOS sean normales.} La
regresión pide normalidad de los \emph{residuos}, no de las variables.
Este Q-Q sirve para conocer la forma de la variable y decidir una
transformación, no para aprobar o reprobar un modelo que todavía no
existe.

\textbf{Con n \textless{} 20 la nube es tan errática que casi nada se
puede concluir.} Se espera ver zigzag; no lo leas como estructura.

Parámetros:

{\def\LTcaptype{none} % do not increment counter
\begin{longtable}[]{@{}ll@{}}
\toprule\noalign{}
Campo & Valor \\
\midrule\noalign{}
\endhead
\bottomrule\noalign{}
\endlastfoot
variable & Temperatura \\
\end{longtable}
}

\begin{Shaded}
\begin{Highlighting}[]
\CommentTok{\# f1.analisis.qq\_normal\_datos}
\NormalTok{x }\OtherTok{\textless{}{-}}\NormalTok{ datos}\SpecialCharTok{$}\NormalTok{Temperatura}
\FunctionTok{qqnorm}\NormalTok{(x, }\AttributeTok{main =} \StringTok{"Q{-}Q normal · Temperatura"}\NormalTok{)}
\FunctionTok{qqline}\NormalTok{(x, }\AttributeTok{col =} \StringTok{"firebrick"}\NormalTok{)}
\end{Highlighting}
\end{Shaded}

\pandocbounded{\includegraphics[keepaspectratio]{taller-01_files/taller-01_files/figure-latex/unnamed-chunk-12-1.pdf}}

\begin{Shaded}
\begin{Highlighting}[]
\FunctionTok{cat}\NormalTok{(}\StringTok{"Shapiro{-}Wilk:}\SpecialCharTok{\textbackslash{}n}\StringTok{"}\NormalTok{); }\FunctionTok{print}\NormalTok{(}\FunctionTok{shapiro.test}\NormalTok{(x))}
\end{Highlighting}
\end{Shaded}

\begin{verbatim}
## Shapiro-Wilk:
\end{verbatim}

\begin{verbatim}
## 
##  Shapiro-Wilk normality test
## 
## data:  x
## W = 0.85961, p-value = 3.403e-09
\end{verbatim}

\begin{quote}
Interpretá acá qué muestra el gráfico y qué conclusión sacás. Una figura
sin lectura no es un resultado: es un adorno.
\end{quote}

\subsection{Histograma · Presion}\label{histograma-presion}

Panel \texttt{f1.analisis.histograma} · marcado a las 07:25:51

El texto de este panel está en la primera sección «Histograma».

Parámetros:

{\def\LTcaptype{none} % do not increment counter
\begin{longtable}[]{@{}ll@{}}
\toprule\noalign{}
Campo & Valor \\
\midrule\noalign{}
\endhead
\bottomrule\noalign{}
\endlastfoot
variable & Presion \\
densidad & TRUE \\
\end{longtable}
}

\begin{Shaded}
\begin{Highlighting}[]
\CommentTok{\# f1.analisis.histograma}
\FunctionTok{ggplot}\NormalTok{(datos, }\FunctionTok{aes}\NormalTok{(}\AttributeTok{x =}\NormalTok{ Presion)) }\SpecialCharTok{+}
  \FunctionTok{geom\_histogram}\NormalTok{(}\FunctionTok{aes}\NormalTok{(}\AttributeTok{y =} \FunctionTok{after\_stat}\NormalTok{(density)),}
                 \AttributeTok{bins =} \DecValTok{30}\NormalTok{, }\AttributeTok{fill =} \StringTok{"steelblue"}\NormalTok{) }\SpecialCharTok{+}
  \FunctionTok{geom\_density}\NormalTok{(}\AttributeTok{alpha =} \FloatTok{0.4}\NormalTok{) }\SpecialCharTok{+}
  \FunctionTok{labs}\NormalTok{(}\AttributeTok{x =} \StringTok{"Presion"}\NormalTok{, }\AttributeTok{y =} \StringTok{"densidad"}\NormalTok{)}
\end{Highlighting}
\end{Shaded}

\pandocbounded{\includegraphics[keepaspectratio]{taller-01_files/taller-01_files/figure-latex/unnamed-chunk-13-1.pdf}}

\begin{quote}
Interpretá acá qué muestra el gráfico y qué conclusión sacás. Una figura
sin lectura no es un resultado: es un adorno.
\end{quote}

\subsection{Densidad kernel ·
Velocidad\_del\_Viento}\label{densidad-kernel-velocidad_del_viento}

Panel \texttt{f1.analisis.densidad} · marcado a las 07:25:51

\subsubsection{Para qué sirve}\label{para-quuxe9-sirve-9}

Ver la forma de una variable sin la arbitrariedad del número de clases.
Es la que conviene para juzgar si una transformación hizo falta y si
funcionó, porque el cambio de forma se lee sin depender de dónde cayeron
los bordes.

\subsubsection{Qué muestra}\label{quuxe9-muestra-9}

Una curva suave que estima la densidad de probabilidad de la variable.
En vez de contar por intervalos, pone una campanita ---el
\textbf{núcleo}--- sobre cada observación y las suma todas.

El parámetro que manda es el ancho de banda \texttt{h}, no el número de
clases:

\[
\hat{f}(x) = \frac{1}{n h} \sum_{i=1}^{n} K\!\left( \frac{x - x_i}{h} \right)
\]

El área bajo la curva vale 1, así que la altura no es una cuenta de
casos: es densidad. Comparar dos densidades de muestras de tamaño
distinto es legítimo; comparar dos histogramas de conteos, no.

\subsubsection{Qué buscar}\label{quuxe9-buscar-9}

\begin{itemize}
\tightlist
\item
  \textbf{Modas}: los picos son la razón de ser de este gráfico. Dos
  modas separadas y estables al mover \texttt{h} son dos poblaciones
  mezcladas.
\item
  \textbf{Asimetría}: la cola más larga dice hacia dónde se estira la
  variable.
\item
  \textbf{Colas}: qué tan rápido cae la curva en los extremos, comparado
  con la campana que uno esperaría si fuera normal.
\item
  \textbf{El valor de h}: aparece en el subtítulo. Por defecto sale de
  la regla de Silverman,
  \texttt{h\ =\ 0.9\ ·\ min(s,\ RIC/1.34)\ ·\ n\^{}(−1/5)}, que es un
  punto de partida razonable y nada más.
\end{itemize}

\subsubsection{Cuándo engaña}\label{cuuxe1ndo-engauxf1a-9}

\textbf{h chico inventa modas; h grande las borra.} Es el mismo dilema
del número de clases del histograma, con otro nombre: sesgo contra
varianza. Movelo y fijate cuáles picos sobreviven. Los que aparecen y
desaparecen no son estructura, son ruido suavizado con distinta fuerza.

\textbf{Suaviza más allá del recorrido real.} El núcleo gaussiano tiene
colas infinitas, así que la curva pone densidad en valores imposibles:
edades negativas, porcentajes por encima de 100. Si la variable está
acotada, la parte que se derrama por el borde es un artefacto del
método.

\textbf{Con pocos datos es casi solo el núcleo.} Con n = 10 estás
mirando la forma de la campanita que elegiste, no la de la variable.

Parámetros:

{\def\LTcaptype{none} % do not increment counter
\begin{longtable}[]{@{}ll@{}}
\toprule\noalign{}
Campo & Valor \\
\midrule\noalign{}
\endhead
\bottomrule\noalign{}
\endlastfoot
variable & Velocidad\_del\_Viento \\
normal & TRUE \\
\end{longtable}
}

\begin{Shaded}
\begin{Highlighting}[]
\CommentTok{\# f1.analisis.densidad}
\FunctionTok{ggplot}\NormalTok{(datos, }\FunctionTok{aes}\NormalTok{(}\AttributeTok{x =}\NormalTok{ Velocidad\_del\_Viento)) }\SpecialCharTok{+}
  \FunctionTok{geom\_density}\NormalTok{() }\SpecialCharTok{+}
  \FunctionTok{stat\_function}\NormalTok{(}\AttributeTok{fun =}\NormalTok{ dnorm, }\AttributeTok{args =} \FunctionTok{list}\NormalTok{(}\AttributeTok{mean =} \FunctionTok{mean}\NormalTok{(datos}\SpecialCharTok{$}\NormalTok{Velocidad\_del\_Viento, }\AttributeTok{na.rm =} \ConstantTok{TRUE}\NormalTok{), }\AttributeTok{sd =} \FunctionTok{sd}\NormalTok{(datos}\SpecialCharTok{$}\NormalTok{Velocidad\_del\_Viento, }\AttributeTok{na.rm =} \ConstantTok{TRUE}\NormalTok{)), }\AttributeTok{color =} \StringTok{"firebrick"}\NormalTok{, }\AttributeTok{linetype =} \StringTok{"dashed"}\NormalTok{) }\SpecialCharTok{+}
  \FunctionTok{labs}\NormalTok{(}\AttributeTok{x =} \StringTok{"Valor"}\NormalTok{, }\AttributeTok{y =} \StringTok{"densidad"}\NormalTok{)}
\end{Highlighting}
\end{Shaded}

\pandocbounded{\includegraphics[keepaspectratio]{taller-01_files/taller-01_files/figure-latex/unnamed-chunk-14-1.pdf}}

\begin{quote}
Interpretá acá qué muestra el gráfico y qué conclusión sacás. Una figura
sin lectura no es un resultado: es un adorno.
\end{quote}

\begin{center}\rule{0.5\linewidth}{0.5pt}\end{center}

14 paneles · dataset ori · generado en modo servidor.

\end{document}
